% Recovered from the previously completed PDF; no new mathematical content.
% Prose is editable; formula boxes retain original symbol positions.
% Build from repository root: xelatex -no-shell-escape -halt-on-error -output-directory=... print/fudan/NAME.tex
\documentclass[11pt,a4paper]{article}
\usepackage[margin=23mm,headheight=14pt,headsep=8.5mm,footskip=10mm]{geometry}
\usepackage{amsmath,amssymb}
\DeclareMathAlphabet{\mathsf}{OT1}{cmss}{m}{n}
\usepackage[no-math]{fontspec}
\usepackage{xeCJK}
\setmainfont{Latin Modern Roman}
\setCJKmainfont{FandolSong-Regular.otf}[Path=assets/fonts/fandol/,BoldFont=FandolSong-Bold.otf,ItalicFont=FandolSong-Regular.otf,Script=CustomDefault]
\xeCJKsetup{PunctStyle=quanjiao,CheckSingle=true}
\usepackage{setspace}\setstretch{1.12}
\setlength{\parindent}{2em}\setlength{\parskip}{3pt}
\emergencystretch=6em
\hfuzz=9pt
\usepackage{titlesec}
\titleformat{\section}{\centering\normalfont\Large\bfseries}{}{0pt}{}
\titlespacing*{\section}{0pt}{1.8\baselineskip}{1.0\baselineskip}
\usepackage{fancyhdr}\pagestyle{fancy}\fancyhf{}
\fancyhead[L]{\small 数学手记·高等代数}\fancyhead[R]{\small 18S03–18S09}
\fancyfoot[C]{\thepage}\renewcommand{\headrulewidth}{0.4pt}
\usepackage[hidelinks,unicode]{hyperref}
\hypersetup{pdftitle={18S03–18S09：题目、解答与小结},pdfauthor={Vincent Arqain}}
\newcommand{\RB}[4]{\raisebox{0pt}[#2bp][#3bp]{\hbox to #1bp{#4\hss}}}
\newcommand{\RG}[5]{\hbox to 0pt{\kern#1bp\raisebox{#2bp}[0pt][0pt]{\fontsize{#3bp}{#3bp}\csname font#4\endcsname\char#5}\hss}}
\newcommand{\RT}[4]{\hbox to 0pt{\kern#1bp\raisebox{#2bp}[0pt][0pt]{\fontsize{#3bp}{#3bp}\selectfont#4}\hss}}
\newcommand{\RS}[4]{\hbox to 0pt{\kern#1bp\raisebox{#2bp}[0pt][0pt]{\fontsize{#3bp}{#3bp}\selectfont\ensuremath{#4}}\hss}}
\newcommand{\RR}[4]{\hbox to 0pt{\kern#1bp\raisebox{#2bp}[0pt][0pt]{\rule{#3bp}{#4bp}}\hss}}
\newcommand{\fontmi}{\usefont{OML}{cmm}{m}{it}}
\newcommand{\fontsy}{\usefont{OMS}{cmsy}{m}{n}}
\newcommand{\fontex}{\usefont{OMX}{cmex}{m}{n}}
\newcommand{\fontrm}{\usefont{OT1}{cmr}{m}{n}}
\newcommand{\fontbf}{\usefont{OT1}{cmr}{bx}{n}}
\newcommand{\fontam}{\usefont{U}{msa}{m}{n}}
\newcommand{\fontbb}{\usefont{U}{msb}{m}{n}}
\widowpenalty=10000\clubpenalty=10000\raggedbottom
\begin{document}
\begin{center}{\LARGE\bfseries 18S03–18S09\par}\vspace{8pt}复旦大学高等代数习题集·题目与解答\end{center}
% Original page 1.
复旦大学高等代数习题集2023版（谢启鸿）

题目选自原书第60–61页。沿用\allowbreak{}\hskip.1em \RB{11.521}{8.182}{0.500}{%
\RG{0.000}{0.004}{10.9091}{mi}{88}\RG{9.028}{3.616}{7.9701}{sy}{48}%
}\hskip.1em\allowbreak{} 表示转置，\allowbreak{}\hskip.1em \RB{11.521}{10.311}{0.500}{%
\RG{0.000}{0.004}{10.9091}{mi}{88}\RG{9.028}{5.745}{7.9701}{sy}{48}\RR{0.000}{8.078}{9.028}{0.364}%
}\hskip.1em\allowbreak{} 表示共轭转置。

\section*{2018S03}
有限维线性空间上的线性变换至多只有有限个特征值。试构造无限维线性空间\allowbreak{}\hskip.1em \RB{8.715}{7.952}{0.728}{%
\RG{0.000}{0.001}{10.9091}{mi}{86}%
}\hskip.1em\allowbreak{} 上的线性变换\allowbreak{}\hskip.1em \RB{7.122}{5.311}{2.856}{%
\RG{0.000}{-0.000}{10.9091}{mi}{39}%
}\hskip.1em\allowbreak{} ，使得\allowbreak{}\hskip.1em \RB{7.122}{5.311}{2.856}{%
\RG{0.000}{-0.000}{10.9091}{mi}{39}%
}\hskip.1em\allowbreak{} 有无限个特征值。

解：取\allowbreak{}\hskip.1em \RB{42.140}{8.680}{3.229}{%
\RG{0.000}{-0.002}{10.9091}{mi}{86}\RG{11.631}{-0.002}{10.9091}{rm}{61}\RG{22.970}{-0.002}{10.9091}{bb}{82}\RG{30.848}{-0.002}{10.9091}{rm}{91}\RG{33.879}{-0.002}{10.9091}{mi}{120}\RG{39.949}{-0.002}{10.9091}{rm}{93}%
}\hskip.1em\allowbreak{} ，定义

\[
\RB{148.517}{9.069}{3.239}{%
\RG{129.080}{4.503}{7.9701}{sy}{48}\RG{0.000}{-0.001}{10.9091}{mi}{39}\RG{8.210}{-0.001}{10.9091}{rm}{58}\RG{14.814}{-0.001}{10.9091}{mi}{86}\RG{26.445}{-0.001}{10.9091}{sy}{0}\RG{33.121}{-0.001}{10.9091}{sy}{33}\RG{47.051}{-0.001}{10.9091}{mi}{86}\RG{53.869}{-0.001}{10.9091}{mi}{59}\RG{80.477}{-0.001}{10.9091}{mi}{39}\RG{87.473}{-0.001}{10.9091}{rm}{40}\RG{91.627}{-0.001}{10.9091}{mi}{102}\RG{98.058}{-0.001}{10.9091}{rm}{41}\RG{105.236}{-0.001}{10.9091}{rm}{61}\RG{116.580}{-0.001}{10.9091}{mi}{120}\RG{122.645}{-0.001}{10.9091}{mi}{102}\RG{131.631}{-0.001}{10.9091}{rm}{40}\RG{135.785}{-0.001}{10.9091}{mi}{120}\RG{141.855}{-0.001}{10.9091}{rm}{41}\RG{146.009}{-0.001}{10.9091}{mi}{58}%
}
\]
由多项式求导的线性性，对任意\allowbreak{}\hskip.1em \RB{36.437}{8.068}{2.608}{%
\RG{0.000}{-0.003}{10.9091}{mi}{97}\RG{5.607}{-0.003}{10.9091}{mi}{59}\RG{10.386}{-0.003}{10.9091}{mi}{98}\RG{17.965}{-0.003}{10.9091}{sy}{50}\RG{28.268}{-0.003}{10.9091}{bb}{82}%
}\hskip.1em\allowbreak{} 及\allowbreak{}\hskip.1em \RB{38.176}{8.177}{2.728}{%
\RG{0.000}{-0.003}{10.9091}{mi}{102}\RG{5.836}{-0.003}{10.9091}{mi}{59}\RG{10.626}{-0.003}{10.9091}{mi}{103}\RG{19.158}{-0.003}{10.9091}{sy}{50}\RG{29.461}{-0.003}{10.9091}{mi}{86}%
}\hskip.1em\allowbreak{} ，有

\[
\RB{133.398}{8.681}{3.239}{%
\RG{0.000}{-0.001}{10.9091}{mi}{39}\RG{6.998}{-0.001}{10.9091}{rm}{40}\RG{11.152}{-0.001}{10.9091}{mi}{97}\RG{16.759}{-0.001}{10.9091}{mi}{102}\RG{25.613}{-0.001}{10.9091}{rm}{43}\RG{36.346}{-0.001}{10.9091}{mi}{98}\RG{40.885}{-0.001}{10.9091}{mi}{103}\RG{46.394}{-0.001}{10.9091}{rm}{41}\RG{53.572}{-0.001}{10.9091}{rm}{61}\RG{64.917}{-0.001}{10.9091}{mi}{97}\RG{70.524}{-0.001}{10.9091}{mi}{39}\RG{77.522}{-0.001}{10.9091}{rm}{40}\RG{81.675}{-0.001}{10.9091}{mi}{102}\RG{88.106}{-0.001}{10.9091}{rm}{41}\RG{94.684}{-0.001}{10.9091}{rm}{43}\RG{105.416}{-0.001}{10.9091}{mi}{98}\RG{109.954}{-0.001}{10.9091}{mi}{39}\RG{116.956}{-0.001}{10.9091}{rm}{40}\RG{121.110}{-0.001}{10.9091}{mi}{103}\RG{126.616}{-0.001}{10.9091}{rm}{41}\RG{130.770}{-0.001}{10.9091}{mi}{59}%
}
\]
故\allowbreak{}\hskip.1em \RB{7.122}{5.312}{2.855}{%
\RG{0.000}{0.001}{10.9091}{mi}{39}%
}\hskip.1em\allowbreak{} 是线性变换。由于\allowbreak{}\hskip.1em \RB{48.233}{9.297}{2.605}{%
\RG{0.000}{0.001}{10.9091}{rm}{49}\RG{5.341}{0.001}{10.9091}{mi}{59}\RG{10.119}{0.001}{10.9091}{mi}{120}\RG{16.185}{0.001}{10.9091}{mi}{59}\RG{20.974}{0.001}{10.9091}{mi}{120}\RG{27.051}{3.960}{7.9701}{rm}{50}\RG{31.369}{0.001}{10.9091}{mi}{59}\RG{36.148}{0.001}{10.9091}{mi}{58}\RG{40.937}{0.001}{10.9091}{mi}{58}\RG{45.726}{0.001}{10.9091}{mi}{58}%
}\hskip.1em\allowbreak{} 线性无关，\allowbreak{}\hskip.1em \RB{8.715}{7.952}{0.728}{%
\RG{0.000}{0.001}{10.9091}{mi}{86}%
}\hskip.1em\allowbreak{} 是无限维线性空间。对每个整数\allowbreak{}\hskip.1em \RB{26.875}{8.072}{1.994}{%
\RG{0.000}{0.001}{10.9091}{mi}{107}\RG{9.463}{0.001}{10.9091}{sy}{21}\RG{21.488}{0.001}{10.9091}{rm}{48}%
}\hskip.1em\allowbreak{} ，有

\[
\RB{59.707}{10.053}{3.237}{%
\RG{17.221}{4.506}{7.9701}{mi}{107}\RG{52.407}{4.506}{7.9701}{mi}{107}\RG{0.000}{0.001}{10.9091}{mi}{39}\RG{6.997}{0.001}{10.9091}{rm}{40}\RG{11.151}{0.001}{10.9091}{mi}{120}\RG{21.893}{0.001}{10.9091}{rm}{41}\RG{29.071}{0.001}{10.9091}{rm}{61}\RG{40.416}{0.001}{10.9091}{mi}{107}\RG{46.329}{0.001}{10.9091}{mi}{120}\RG{57.079}{0.001}{10.9091}{mi}{59}%
}
\]
且\allowbreak{}\hskip.1em \RB{30.495}{9.506}{2.856}{%
\RG{0.000}{0.000}{10.9091}{mi}{120}\RG{6.070}{3.959}{7.9701}{mi}{107}\RG{13.773}{0.000}{10.9091}{sy}{54}\RG{13.773}{0.000}{10.9091}{rm}{61}\RG{25.107}{0.000}{10.9091}{rm}{48}%
}\hskip.1em\allowbreak{} ，所以\allowbreak{}\hskip.1em \RB{42.457}{7.755}{2.605}{%
\RG{0.000}{0.000}{10.9091}{rm}{48}\RG{5.341}{0.000}{10.9091}{mi}{59}\RG{10.126}{0.000}{10.9091}{rm}{49}\RG{15.467}{0.000}{10.9091}{mi}{59}\RG{20.252}{0.000}{10.9091}{rm}{50}\RG{25.593}{0.000}{10.9091}{mi}{59}\RG{30.372}{0.000}{10.9091}{mi}{58}\RG{35.161}{0.000}{10.9091}{mi}{58}\RG{39.950}{0.000}{10.9091}{mi}{58}%
}\hskip.1em\allowbreak{} 都是\allowbreak{}\hskip.1em \RB{7.122}{5.311}{2.856}{%
\RG{0.000}{0.000}{10.9091}{mi}{39}%
}\hskip.1em\allowbreak{} 的特征值。

小结：在无限维空间的一组基上指定两两不同的伸缩系数，即可构造具有无限个特征值的线性变换。

\section*{2018S04}
设\allowbreak{}\hskip.1em \RB{8.202}{8.293}{0.500}{%
\RG{0.000}{0.004}{10.9091}{mi}{65}%
}\hskip.1em\allowbreak{} 为\allowbreak{}\hskip.1em \RB{6.555}{5.315}{0.605}{%
\RG{0.000}{0.004}{10.9091}{mi}{110}%
}\hskip.1em\allowbreak{} 阶复矩阵，\allowbreak{}\hskip.1em \RB{18.315}{8.195}{2.613}{%
\RG{0.000}{0.004}{10.9091}{mi}{11}\RG{6.862}{0.004}{10.9091}{mi}{59}\RG{11.651}{0.004}{10.9091}{mi}{12}%
}\hskip.1em\allowbreak{} 为\allowbreak{}\hskip.1em \RB{6.555}{5.315}{0.605}{%
\RG{0.000}{0.004}{10.9091}{mi}{110}%
}\hskip.1em\allowbreak{} 维复列向量，\allowbreak{}\hskip.1em \RB{47.588}{8.528}{2.613}{%
\RG{0.000}{0.004}{10.9091}{mi}{66}\RG{11.973}{0.004}{10.9091}{rm}{61}\RG{23.588}{0.004}{10.9091}{mi}{65}\RG{31.596}{0.004}{10.9091}{mi}{11}\RG{38.458}{0.004}{10.9091}{mi}{12}\RG{45.095}{3.962}{7.9701}{sy}{48}%
}\hskip.1em\allowbreak{} 。试求矩阵\allowbreak{}\hskip.1em \RB{8.606}{7.955}{0.500}{%
\RG{0.000}{0.004}{10.9091}{mi}{66}%
}\hskip.1em\allowbreak{} 可对角化的充要条件。

解：所求充要条件为

\allowbreak{}\hskip.1em \RB{31.455}{8.182}{0.727}{%
\RG{0.000}{0.002}{10.9091}{mi}{66}\RG{11.696}{0.002}{10.9091}{rm}{61}\RG{23.035}{0.002}{10.9091}{mi}{79}%
}\hskip.1em\allowbreak{} 或\allowbreak{}\hskip.1em \RB{46.394}{9.072}{2.855}{%
\RG{0.000}{0.002}{10.9091}{mi}{12}\RG{6.636}{4.507}{7.9701}{sy}{48}\RG{9.186}{0.002}{10.9091}{mi}{65}\RG{17.194}{0.002}{10.9091}{mi}{11}\RG{27.088}{0.002}{10.9091}{sy}{54}\RG{27.088}{0.002}{10.9091}{rm}{61}\RG{38.423}{0.002}{10.9091}{rm}{48}\RG{43.767}{0.002}{10.9091}{mi}{59}%
}\hskip.1em\allowbreak{}亦即\allowbreak{}\hskip.1em \RB{34.623}{8.290}{0.728}{%
\RG{0.000}{0.001}{10.9091}{mi}{65}\RG{8.007}{0.001}{10.9091}{mi}{11}\RG{17.901}{0.001}{10.9091}{rm}{61}\RG{29.236}{0.001}{10.9091}{rm}{48}%
}\hskip.1em\allowbreak{} 、\allowbreak{}\hskip.1em \RB{26.387}{8.192}{2.615}{%
\RG{0.000}{0.001}{10.9091}{mi}{12}\RG{9.665}{0.001}{10.9091}{rm}{61}\RG{21.000}{0.001}{10.9091}{rm}{48}%
}\hskip.1em\allowbreak{} 、\allowbreak{}\hskip.1em \RB{43.810}{8.525}{2.855}{%
\RG{0.000}{0.001}{10.9091}{mi}{12}\RG{6.636}{3.959}{7.9701}{sy}{48}\RG{9.186}{0.001}{10.9091}{mi}{65}\RG{17.194}{0.001}{10.9091}{mi}{11}\RG{27.088}{0.001}{10.9091}{sy}{54}\RG{27.088}{0.001}{10.9091}{rm}{61}\RG{38.423}{0.001}{10.9091}{rm}{48}%
}\hskip.1em\allowbreak{} 三者至少有一个成立。

令\allowbreak{}\hskip.1em \RB{35.381}{8.286}{0.612}{%
\RG{0.000}{-0.003}{10.9091}{mi}{117}\RG{9.110}{-0.003}{10.9091}{rm}{61}\RG{20.448}{-0.003}{10.9091}{mi}{65}\RG{28.456}{-0.003}{10.9091}{mi}{11}%
}\hskip.1em\allowbreak{} ，\allowbreak{}\hskip.1em \RB{34.412}{8.522}{2.619}{%
\RG{0.000}{-0.003}{10.9091}{mi}{99}\RG{7.628}{-0.003}{10.9091}{rm}{61}\RG{18.966}{-0.003}{10.9091}{mi}{12}\RG{25.601}{3.957}{7.9701}{sy}{48}\RG{28.152}{-0.003}{10.9091}{mi}{117}%
}\hskip.1em\allowbreak{} ，则

\[
\RB{161.226}{9.841}{3.239}{%
\RG{35.749}{4.504}{7.9701}{sy}{48}\RG{73.569}{4.504}{7.9701}{rm}{50}\RG{109.125}{4.504}{7.9701}{sy}{48}\RG{128.544}{4.504}{7.9701}{sy}{48}\RG{0.000}{-0.001}{10.9091}{mi}{66}\RG{11.696}{-0.001}{10.9091}{rm}{61}\RG{23.035}{-0.001}{10.9091}{mi}{117}\RG{29.111}{-0.001}{10.9091}{mi}{12}\RG{38.300}{-0.001}{10.9091}{mi}{59}\RG{64.897}{-0.001}{10.9091}{mi}{66}\RG{80.917}{-0.001}{10.9091}{rm}{61}\RG{92.256}{-0.001}{10.9091}{mi}{117}\RG{98.336}{-0.001}{10.9091}{rm}{40}\RG{102.490}{-0.001}{10.9091}{mi}{12}\RG{111.675}{-0.001}{10.9091}{mi}{117}\RG{117.755}{-0.001}{10.9091}{rm}{41}\RG{121.909}{-0.001}{10.9091}{mi}{12}\RG{134.125}{-0.001}{10.9091}{rm}{61}\RG{145.464}{-0.001}{10.9091}{mi}{99}\RG{150.056}{-0.001}{10.9091}{mi}{66}\RG{158.718}{-0.001}{10.9091}{mi}{58}%
}
\]
若\allowbreak{}\hskip.1em \RB{31.455}{8.181}{0.728}{%
\RG{0.000}{0.001}{10.9091}{mi}{66}\RG{11.696}{0.001}{10.9091}{rm}{61}\RG{23.035}{0.001}{10.9091}{mi}{79}%
}\hskip.1em\allowbreak{} ，则\allowbreak{}\hskip.1em \RB{8.606}{7.952}{0.500}{%
\RG{0.000}{0.001}{10.9091}{mi}{66}%
}\hskip.1em\allowbreak{} 可对角化。若\allowbreak{}\hskip.1em \RB{24.350}{8.312}{2.855}{%
\RG{0.000}{0.001}{10.9091}{mi}{99}\RG{7.628}{0.001}{10.9091}{sy}{54}\RG{7.628}{0.001}{10.9091}{rm}{61}\RG{18.962}{0.001}{10.9091}{rm}{48}%
}\hskip.1em\allowbreak{} ，则\allowbreak{}\hskip.1em \RB{25.833}{8.312}{2.855}{%
\RG{0.000}{0.001}{10.9091}{mi}{117}\RG{9.111}{0.001}{10.9091}{sy}{54}\RG{9.111}{0.001}{10.9091}{rm}{61}\RG{20.445}{0.001}{10.9091}{rm}{48}%
}\hskip.1em\allowbreak{} ，且任意\allowbreak{}\hskip.1em \RB{32.233}{8.181}{0.935}{%
\RG{0.000}{0.001}{10.9091}{mi}{120}\RG{9.100}{0.001}{10.9091}{sy}{50}\RG{19.403}{0.001}{10.9091}{bb}{67}\RG{27.282}{3.960}{7.9701}{mi}{110}%
}\hskip.1em\allowbreak{} 均可写为

\[
\RB{122.510}{16.867}{11.397}{%
\RG{58.931}{16.004}{9.0910}{ex}{32}\RG{110.867}{16.004}{9.0910}{ex}{33}\RG{28.165}{11.343}{7.9701}{sy}{48}\RG{93.258}{11.343}{7.9701}{sy}{48}\RG{21.530}{7.383}{10.9091}{mi}{12}\RG{30.715}{7.383}{10.9091}{mi}{120}\RG{86.623}{7.383}{10.9091}{mi}{12}\RG{95.808}{7.383}{10.9091}{mi}{120}\RG{0.000}{0.003}{10.9091}{mi}{120}\RG{9.100}{0.003}{10.9091}{rm}{61}\RG{39.694}{0.003}{10.9091}{mi}{117}\RG{48.199}{0.003}{10.9091}{rm}{43}\RG{66.128}{0.003}{10.9091}{mi}{120}\RG{74.622}{0.003}{10.9091}{sy}{0}\RG{104.787}{0.003}{10.9091}{mi}{117}\RG{119.882}{0.003}{10.9091}{mi}{59}\RG{26.859}{-7.480}{10.9091}{mi}{99}\RG{91.952}{-7.480}{10.9091}{mi}{99}\RR{21.530}{2.730}{15.255}{0.364}\RR{86.623}{2.730}{15.255}{0.364}%
}
\]
其中括号内的向量属于\allowbreak{}\hskip.1em \RB{25.189}{8.520}{2.620}{%
\RG{0.000}{-0.004}{10.9091}{rm}{107}\RG{5.346}{-0.004}{10.9091}{rm}{101}\RG{10.091}{-0.004}{10.9091}{rm}{114}\RG{16.060}{-0.004}{10.9091}{mi}{12}\RG{22.696}{3.955}{7.9701}{sy}{48}%
}\hskip.1em\allowbreak{} 。又\allowbreak{}\hskip.1em \RB{53.985}{8.520}{2.860}{%
\RG{0.000}{-0.004}{10.9091}{mi}{12}\RG{6.636}{3.955}{7.9701}{sy}{48}\RG{9.186}{-0.004}{10.9091}{mi}{117}\RG{18.296}{-0.004}{10.9091}{rm}{61}\RG{29.635}{-0.004}{10.9091}{mi}{99}\RG{37.262}{-0.004}{10.9091}{sy}{54}\RG{37.262}{-0.004}{10.9091}{rm}{61}\RG{48.597}{-0.004}{10.9091}{rm}{48}%
}\hskip.1em\allowbreak{} ，故

\[
\RB{226.052}{11.598}{7.822}{%
\RG{181.593}{10.907}{9.0910}{ex}{12}\RG{181.593}{5.452}{9.0910}{ex}{12}\RG{7.879}{4.502}{7.9701}{mi}{110}\RG{77.378}{4.502}{7.9701}{sy}{48}\RG{0.000}{-0.002}{10.9091}{bb}{67}\RG{16.052}{-0.002}{10.9091}{rm}{61}\RG{27.391}{-0.002}{10.9091}{bb}{67}\RG{35.270}{-0.002}{10.9091}{mi}{117}\RG{43.773}{-0.002}{10.9091}{sy}{8}\RG{54.682}{-0.002}{10.9091}{rm}{107}\RG{60.028}{-0.002}{10.9091}{rm}{101}\RG{64.773}{-0.002}{10.9091}{rm}{114}\RG{70.743}{-0.002}{10.9091}{mi}{12}\RG{79.929}{-0.002}{10.9091}{mi}{59}\RG{106.526}{-0.002}{10.9091}{mi}{66}\RG{115.198}{-0.002}{10.9091}{mi}{117}\RG{124.308}{-0.002}{10.9091}{rm}{61}\RG{135.647}{-0.002}{10.9091}{mi}{99}\RG{140.240}{-0.002}{10.9091}{mi}{117}\RG{146.316}{-0.002}{10.9091}{mi}{59}\RG{172.913}{-0.002}{10.9091}{mi}{66}\RG{181.593}{-0.002}{9.0910}{ex}{12}\RG{206.865}{-0.002}{10.9091}{rm}{61}\RG{218.200}{-0.002}{10.9091}{rm}{48}\RG{223.545}{-0.002}{10.9091}{mi}{58}\RG{200.930}{-3.841}{5.9776}{sy}{48}\RG{184.623}{-5.911}{7.9701}{rm}{107}\RG{188.492}{-5.911}{7.9701}{rm}{101}\RG{191.925}{-5.911}{7.9701}{rm}{114}\RG{196.220}{-5.911}{7.9701}{mi}{12}%
}
\]
选取与此直和分解相适应的基，即得\allowbreak{}\hskip.1em \RB{94.665}{8.684}{3.236}{%
\RG{0.000}{0.002}{10.9091}{mi}{66}\RG{11.696}{0.002}{10.9091}{sy}{24}\RG{23.211}{0.002}{10.9091}{rm}{100}\RG{29.145}{0.002}{10.9091}{rm}{105}\RG{32.113}{0.002}{10.9091}{rm}{97}\RG{37.458}{0.002}{10.9091}{rm}{103}\RG{42.946}{0.002}{10.9091}{rm}{40}\RG{47.093}{0.002}{10.9091}{mi}{99}\RG{51.686}{0.002}{10.9091}{mi}{59}\RG{56.475}{0.002}{10.9091}{rm}{48}\RG{61.816}{0.002}{10.9091}{mi}{59}\RG{66.595}{0.002}{10.9091}{mi}{58}\RG{71.384}{0.002}{10.9091}{mi}{58}\RG{76.173}{0.002}{10.9091}{mi}{58}\RG{80.951}{0.002}{10.9091}{mi}{59}\RG{85.743}{0.002}{10.9091}{rm}{48}\RG{91.089}{0.002}{10.9091}{rm}{41}%
}\hskip.1em\allowbreak{} 。

反之，若\allowbreak{}\hskip.1em \RB{31.844}{8.310}{2.857}{%
\RG{0.000}{-0.001}{10.9091}{mi}{66}\RG{11.891}{-0.001}{10.9091}{sy}{54}\RG{11.891}{-0.001}{10.9091}{rm}{61}\RG{23.424}{-0.001}{10.9091}{mi}{79}%
}\hskip.1em\allowbreak{} 且\allowbreak{}\hskip.1em \RB{24.741}{7.754}{0.730}{%
\RG{0.000}{-0.001}{10.9091}{mi}{99}\RG{7.823}{-0.001}{10.9091}{rm}{61}\RG{19.354}{-0.001}{10.9091}{rm}{48}%
}\hskip.1em\allowbreak{} ，则\allowbreak{}\hskip.1em \RB{36.162}{9.295}{0.730}{%
\RG{0.000}{-0.001}{10.9091}{mi}{66}\RG{8.665}{3.958}{7.9701}{rm}{50}\RG{16.209}{-0.001}{10.9091}{rm}{61}\RG{27.742}{-0.001}{10.9091}{mi}{79}%
}\hskip.1em\allowbreak{} 。设\allowbreak{}\hskip.1em \RB{8.606}{7.950}{0.501}{%
\RG{0.000}{-0.001}{10.9091}{mi}{66}%
}\hskip.1em\allowbreak{} 可对角化，则与它相似的对角阵\allowbreak{}\hskip.1em \RB{9.107}{7.950}{0.501}{%
\RG{0.000}{-0.001}{10.9091}{mi}{68}%
}\hskip.1em\allowbreak{} 满足\allowbreak{}\hskip.1em \RB{36.274}{9.294}{0.730}{%
\RG{0.000}{-0.001}{10.9091}{mi}{68}\RG{9.166}{3.958}{7.9701}{rm}{50}\RG{16.515}{-0.001}{10.9091}{rm}{61}\RG{27.854}{-0.001}{10.9091}{mi}{79}%
}\hskip.1em\allowbreak{} ，从而\allowbreak{}\hskip.1em \RB{31.955}{8.179}{0.730}{%
\RG{0.000}{-0.001}{10.9091}{mi}{68}\RG{12.196}{-0.001}{10.9091}{rm}{61}\RG{23.535}{-0.001}{10.9091}{mi}{79}%
}\hskip.1em\allowbreak{} ，这将推出\allowbreak{}\hskip.1em \RB{31.455}{8.179}{0.730}{%
\RG{0.000}{-0.001}{10.9091}{mi}{66}\RG{11.696}{-0.001}{10.9091}{rm}{61}\RG{23.035}{-0.001}{10.9091}{mi}{79}%
}\hskip.1em\allowbreak{} ，矛盾。小结：秩至多为\allowbreak{}\hskip.1em \RB{4.973}{7.752}{0.503}{%
\RG{0.000}{-0.003}{10.9091}{rm}{49}%
}\hskip.1em\allowbreak{} 的矩阵满足\allowbreak{}\hskip.1em \RB{63.058}{9.292}{3.241}{%
\RG{0.000}{-0.003}{10.9091}{mi}{66}\RG{8.665}{3.956}{7.9701}{rm}{50}\RG{16.014}{-0.003}{10.9091}{rm}{61}\RG{27.348}{-0.003}{10.9091}{rm}{40}\RG{31.505}{-0.003}{10.9091}{rm}{116}\RG{35.661}{-0.003}{10.9091}{rm}{114}\RG{41.632}{-0.003}{10.9091}{mi}{66}\RG{50.299}{-0.003}{10.9091}{rm}{41}\RG{54.452}{-0.003}{10.9091}{mi}{66}%
}\hskip.1em\allowbreak{} ；非零时，可对角化当且仅当其迹不为零。

% Original page 2.
\section*{2018S05}
设\allowbreak{}\hskip.1em \RB{8.649}{8.181}{0.728}{%
\RG{0.000}{0.001}{10.9091}{mi}{67}%
}\hskip.1em\allowbreak{} 为数域\allowbreak{}\hskip.1em \RB{10.035}{7.952}{0.500}{%
\RG{0.000}{0.001}{10.9091}{mi}{75}%
}\hskip.1em\allowbreak{} 上的\allowbreak{}\hskip.1em \RB{6.555}{5.312}{0.608}{%
\RG{0.000}{0.001}{10.9091}{mi}{110}%
}\hskip.1em\allowbreak{} 阶方阵，证明：(1)\allowbreak{}\hskip.1em \RB{8.649}{8.179}{0.730}{%
\RG{0.000}{-0.001}{10.9091}{mi}{67}%
}\hskip.1em\allowbreak{} 是幂零阵当且仅当\allowbreak{}\hskip.1em \RB{8.649}{8.179}{0.730}{%
\RG{0.000}{-0.001}{10.9091}{mi}{67}%
}\hskip.1em\allowbreak{} 的特征值全为零；(2)若\allowbreak{}\hskip.1em \RB{41.188}{8.678}{3.242}{%
\RG{0.000}{-0.004}{10.9091}{mi}{114}\RG{5.111}{-0.004}{10.9091}{rm}{40}\RG{9.265}{-0.004}{10.9091}{mi}{67}\RG{17.691}{-0.004}{10.9091}{rm}{41}\RG{24.869}{-0.004}{10.9091}{rm}{61}\RG{36.215}{-0.004}{10.9091}{rm}{49}%
}\hskip.1em\allowbreak{} ，则\allowbreak{}\hskip.1em \RB{8.649}{8.176}{0.733}{%
\RG{0.000}{-0.004}{10.9091}{mi}{67}%
}\hskip.1em\allowbreak{} 是幂零阵当且仅当\allowbreak{}\hskip.1em \RB{44.799}{8.678}{3.242}{%
\RG{0.000}{-0.004}{10.9091}{rm}{116}\RG{4.156}{-0.004}{10.9091}{rm}{114}\RG{8.313}{-0.004}{10.9091}{rm}{40}\RG{12.462}{-0.004}{10.9091}{mi}{67}\RG{20.887}{-0.004}{10.9091}{rm}{41}\RG{28.066}{-0.004}{10.9091}{rm}{61}\RG{39.411}{-0.004}{10.9091}{rm}{48}%
}\hskip.1em\allowbreak{} ；(3)若\allowbreak{}\hskip.1em \RB{44.799}{8.685}{3.235}{%
\RG{0.000}{0.003}{10.9091}{rm}{116}\RG{4.156}{0.003}{10.9091}{rm}{114}\RG{8.313}{0.003}{10.9091}{rm}{40}\RG{12.462}{0.003}{10.9091}{mi}{67}\RG{20.887}{0.003}{10.9091}{rm}{41}\RG{28.066}{0.003}{10.9091}{rm}{61}\RG{39.411}{0.003}{10.9091}{rm}{48}%
}\hskip.1em\allowbreak{} ，且存在数域\allowbreak{}\hskip.1em \RB{10.035}{7.954}{0.500}{%
\RG{0.000}{0.003}{10.9091}{mi}{75}%
}\hskip.1em\allowbreak{} 上的\allowbreak{}\hskip.1em \RB{6.555}{5.314}{0.606}{%
\RG{0.000}{0.003}{10.9091}{mi}{110}%
}\hskip.1em\allowbreak{} 阶方阵\allowbreak{}\hskip.1em \RB{8.202}{8.292}{0.500}{%
\RG{0.000}{0.003}{10.9091}{mi}{65}%
}\hskip.1em\allowbreak{} ，使得\allowbreak{}\hskip.1em \RB{8.202}{8.292}{0.500}{%
\RG{0.000}{0.003}{10.9091}{mi}{65}%
}\hskip.1em\allowbreak{} 的特征多项式是\allowbreak{}\hskip.1em \RB{10.035}{7.954}{0.500}{%
\RG{0.000}{0.003}{10.9091}{mi}{75}%
}\hskip.1em\allowbreak{} 上的不可

约多项式，以及

\[
\RB{151.725}{9.822}{3.235}{%
\RG{28.896}{4.508}{7.9701}{mi}{105}\RG{0.000}{0.003}{10.9091}{rm}{116}\RG{4.156}{0.003}{10.9091}{rm}{114}\RG{8.313}{0.003}{10.9091}{rm}{40}\RG{12.462}{0.004}{10.9091}{mi}{67}\RG{20.884}{0.004}{10.9091}{mi}{65}\RG{31.981}{0.003}{10.9091}{rm}{41}\RG{39.159}{0.003}{10.9091}{rm}{61}\RG{50.505}{0.003}{10.9091}{rm}{48}\RG{77.669}{0.003}{10.9091}{rm}{40}\RG{81.825}{0.003}{10.9091}{rm}{49}\RG{90.189}{0.004}{10.9091}{sy}{20}\RG{101.704}{0.003}{10.9091}{mi}{105}\RG{108.378}{0.004}{10.9091}{sy}{20}\RG{119.894}{0.003}{10.9091}{mi}{110}\RG{128.694}{0.004}{10.9091}{sy}{0}\RG{139.603}{0.003}{10.9091}{rm}{49}\RG{144.949}{0.003}{10.9091}{rm}{41}\RG{149.098}{0.004}{10.9091}{mi}{59}%
}
\]
则必有\allowbreak{}\hskip.1em \RB{41.182}{8.684}{3.236}{%
\RG{0.000}{0.002}{10.9091}{mi}{114}\RG{5.110}{0.002}{10.9091}{rm}{40}\RG{9.264}{0.002}{10.9091}{mi}{67}\RG{17.690}{0.002}{10.9091}{rm}{41}\RG{24.874}{0.002}{10.9091}{sy}{54}\RG{24.874}{0.002}{10.9091}{rm}{61}\RG{36.209}{0.002}{10.9091}{rm}{49}%
}\hskip.1em\allowbreak{} 。

证明：(1)特征值在\allowbreak{}\hskip.1em \RB{7.973}{8.181}{0.706}{%
\RG{0.000}{0.001}{10.9091}{bb}{67}%
}\hskip.1em\allowbreak{} 中考虑。若\allowbreak{}\hskip.1em \RB{35.380}{8.181}{0.728}{%
\RG{0.000}{0.001}{10.9091}{mi}{67}\RG{8.425}{3.960}{7.9701}{mi}{113}\RG{15.621}{0.001}{10.9091}{rm}{61}\RG{26.960}{0.001}{10.9091}{mi}{79}%
}\hskip.1em\allowbreak{} ，\allowbreak{}\hskip.1em \RB{40.065}{8.181}{0.728}{%
\RG{0.000}{0.001}{10.9091}{mi}{67}\RG{8.422}{0.001}{10.9091}{mi}{118}\RG{17.012}{0.001}{10.9091}{rm}{61}\RG{28.350}{0.001}{10.9091}{mi}{21}\RG{34.579}{0.001}{10.9091}{mi}{118}%
}\hskip.1em\allowbreak{} 且\allowbreak{}\hskip.1em \RB{25.308}{8.312}{2.855}{%
\RG{0.000}{0.001}{10.9091}{mi}{118}\RG{8.586}{0.001}{10.9091}{sy}{54}\RG{8.586}{0.001}{10.9091}{rm}{61}\RG{19.921}{0.001}{10.9091}{rm}{48}%
}\hskip.1em\allowbreak{} ，则

\[
\RB{70.791}{8.214}{2.603}{%
\RG{28.136}{4.507}{7.9701}{mi}{113}\RG{58.456}{4.507}{7.9701}{mi}{113}\RG{0.000}{0.003}{10.9091}{rm}{48}\RG{8.367}{0.003}{10.9091}{rm}{61}\RG{19.709}{0.003}{10.9091}{mi}{67}\RG{32.301}{0.003}{10.9091}{mi}{118}\RG{40.887}{0.003}{10.9091}{rm}{61}\RG{52.226}{0.003}{10.9091}{mi}{21}\RG{62.622}{0.003}{10.9091}{mi}{118}\RG{68.164}{0.003}{10.9091}{mi}{59}%
}
\]
从而\allowbreak{}\hskip.1em \RB{25.983}{8.066}{0.734}{%
\RG{0.000}{-0.005}{10.9091}{mi}{21}\RG{9.261}{-0.005}{10.9091}{rm}{61}\RG{20.596}{-0.005}{10.9091}{rm}{48}%
}\hskip.1em\allowbreak{} 。反之，若\allowbreak{}\hskip.1em \RB{8.649}{8.176}{0.734}{%
\RG{0.000}{-0.005}{10.9091}{mi}{67}%
}\hskip.1em\allowbreak{} 的特征值全为零，则其特征多项式为\allowbreak{}\hskip.1em \RB{8.809}{7.661}{0.614}{%
\RG{0.000}{-0.005}{10.9091}{mi}{116}\RG{3.858}{3.954}{7.9701}{mi}{110}%
}\hskip.1em\allowbreak{} 。由Cayley–Hamilton定理，有\allowbreak{}\hskip.1em \RB{36.358}{8.184}{0.725}{%
\RG{0.000}{0.004}{10.9091}{mi}{67}\RG{8.425}{3.963}{7.9701}{mi}{110}\RG{16.599}{0.004}{10.9091}{rm}{61}\RG{27.937}{0.004}{10.9091}{mi}{79}%
}\hskip.1em\allowbreak{} 。

(2)由\allowbreak{}\hskip.1em \RB{41.187}{8.683}{3.237}{%
\RG{0.000}{0.001}{10.9091}{mi}{114}\RG{5.110}{0.001}{10.9091}{rm}{40}\RG{9.265}{0.001}{10.9091}{mi}{67}\RG{17.690}{0.001}{10.9091}{rm}{41}\RG{24.868}{0.001}{10.9091}{rm}{61}\RG{36.214}{0.001}{10.9091}{rm}{49}%
}\hskip.1em\allowbreak{} ，取\allowbreak{}\hskip.1em \RB{8.649}{8.181}{0.728}{%
\RG{0.000}{0.001}{10.9091}{mi}{67}%
}\hskip.1em\allowbreak{} 的一列非零向量\allowbreak{}\hskip.1em \RB{6.260}{5.312}{0.608}{%
\RG{0.000}{0.001}{10.9091}{mi}{117}%
}\hskip.1em\allowbreak{} ，则每列均为\allowbreak{}\hskip.1em \RB{6.260}{5.312}{0.608}{%
\RG{0.000}{0.001}{10.9091}{mi}{117}%
}\hskip.1em\allowbreak{} 的倍数，故可写成\allowbreak{}\hskip.1em \RB{36.923}{8.526}{0.728}{%
\RG{0.000}{0.001}{10.9091}{mi}{67}\RG{11.456}{0.001}{10.9091}{rm}{61}\RG{22.795}{0.001}{10.9091}{mi}{117}\RG{28.871}{0.001}{10.9091}{mi}{118}\RG{34.430}{3.960}{7.9701}{sy}{48}%
}\hskip.1em\allowbreak{} ，其中\allowbreak{}\hskip.1em \RB{44.576}{7.952}{2.605}{%
\RG{0.000}{0.001}{10.9091}{mi}{117}\RG{6.076}{0.001}{10.9091}{mi}{59}\RG{10.855}{0.001}{10.9091}{mi}{118}\RG{19.450}{0.001}{10.9091}{sy}{50}\RG{29.753}{0.001}{10.9091}{mi}{75}\RG{39.625}{3.959}{7.9701}{mi}{110}%
}\hskip.1em\allowbreak{} 均非零。令

\[
\RB{123.916}{15.529}{12.850}{%
\RG{51.885}{11.822}{7.9701}{mi}{110}\RG{47.664}{9.094}{9.0910}{ex}{88}\RG{24.522}{4.507}{7.9701}{sy}{48}\RG{0.000}{0.003}{10.9091}{mi}{99}\RG{7.628}{0.003}{10.9091}{rm}{61}\RG{18.967}{0.003}{10.9091}{mi}{118}\RG{27.072}{0.003}{10.9091}{mi}{117}\RG{36.183}{0.003}{10.9091}{rm}{61}\RG{62.757}{0.003}{10.9091}{mi}{117}\RG{72.836}{0.003}{10.9091}{mi}{118}\RG{85.029}{0.003}{10.9091}{rm}{61}\RG{96.364}{0.003}{10.9091}{rm}{116}\RG{100.520}{0.003}{10.9091}{rm}{114}\RG{104.676}{0.003}{10.9091}{rm}{40}\RG{108.830}{0.003}{10.9091}{mi}{67}\RG{117.255}{0.003}{10.9091}{rm}{41}\RG{121.409}{0.003}{10.9091}{mi}{58}\RG{68.837}{-1.633}{7.9701}{mi}{106}\RG{78.000}{-1.633}{7.9701}{mi}{106}\RG{47.521}{-10.866}{7.9701}{mi}{106}\RG{51.065}{-10.866}{7.9701}{rm}{61}\RG{57.072}{-10.866}{7.9701}{rm}{49}%
}
\]
直接相乘并归纳得

\[
\RB{166.907}{9.840}{3.239}{%
\RG{8.425}{4.503}{7.9701}{rm}{50}\RG{74.572}{4.503}{7.9701}{mi}{113}\RG{97.704}{4.503}{7.9701}{mi}{113}\RG{101.415}{4.503}{7.9701}{sy}{0}\RG{107.425}{4.503}{7.9701}{rm}{49}\RG{0.000}{-0.001}{10.9091}{mi}{67}\RG{15.774}{-0.001}{10.9091}{rm}{61}\RG{27.113}{-0.001}{10.9091}{mi}{99}\RG{31.705}{-0.001}{10.9091}{mi}{67}\RG{39.538}{-0.001}{10.9091}{mi}{59}\RG{66.146}{-0.001}{10.9091}{mi}{67}\RG{81.768}{-0.001}{10.9091}{rm}{61}\RG{93.106}{-0.001}{10.9091}{mi}{99}\RG{111.743}{-0.001}{10.9091}{mi}{67}\RG{131.078}{-0.001}{10.9091}{rm}{40}\RG{135.232}{-0.001}{10.9091}{mi}{113}\RG{143.389}{-0.001}{10.9091}{sy}{21}\RG{154.905}{-0.001}{10.9091}{rm}{49}\RG{160.251}{-0.001}{10.9091}{rm}{41}\RG{164.400}{-0.001}{10.9091}{mi}{58}%
}
\]
若\allowbreak{}\hskip.1em \RB{24.349}{7.752}{0.732}{%
\RG{0.000}{-0.003}{10.9091}{mi}{99}\RG{7.627}{-0.003}{10.9091}{rm}{61}\RG{18.962}{-0.003}{10.9091}{rm}{48}%
}\hskip.1em\allowbreak{} ，则\allowbreak{}\hskip.1em \RB{35.533}{9.293}{0.732}{%
\RG{0.000}{-0.003}{10.9091}{mi}{67}\RG{8.426}{3.957}{7.9701}{rm}{50}\RG{15.775}{-0.003}{10.9091}{rm}{61}\RG{27.113}{-0.003}{10.9091}{mi}{79}%
}\hskip.1em\allowbreak{} ；若\allowbreak{}\hskip.1em \RB{8.649}{8.177}{0.732}{%
\RG{0.000}{-0.003}{10.9091}{mi}{67}%
}\hskip.1em\allowbreak{} 幂零，取\allowbreak{}\hskip.1em \RB{24.973}{7.752}{2.619}{%
\RG{0.000}{-0.003}{10.9091}{mi}{113}\RG{8.158}{-0.003}{10.9091}{sy}{21}\RG{19.673}{-0.003}{10.9091}{rm}{50}%
}\hskip.1em\allowbreak{} 使\allowbreak{}\hskip.1em \RB{35.380}{8.177}{0.732}{%
\RG{0.000}{-0.003}{10.9091}{mi}{67}\RG{8.425}{3.957}{7.9701}{mi}{113}\RG{15.621}{-0.003}{10.9091}{rm}{61}\RG{26.960}{-0.003}{10.9091}{mi}{79}%
}\hskip.1em\allowbreak{} ，由\allowbreak{}\hskip.1em \RB{31.215}{8.308}{2.859}{%
\RG{0.000}{-0.003}{10.9091}{mi}{67}\RG{11.456}{-0.003}{10.9091}{sy}{54}\RG{11.456}{-0.003}{10.9091}{rm}{61}\RG{22.795}{-0.003}{10.9091}{mi}{79}%
}\hskip.1em\allowbreak{} 得\allowbreak{}\hskip.1em \RB{38.389}{9.293}{0.732}{%
\RG{0.000}{-0.003}{10.9091}{mi}{99}\RG{4.597}{3.957}{7.9701}{mi}{113}\RG{8.308}{3.957}{7.9701}{sy}{0}\RG{14.318}{3.957}{7.9701}{rm}{49}\RG{21.667}{-0.003}{10.9091}{rm}{61}\RG{33.001}{-0.003}{10.9091}{rm}{48}%
}\hskip.1em\allowbreak{} ，故\allowbreak{}\hskip.1em \RB{24.350}{7.752}{0.732}{%
\RG{0.000}{-0.003}{10.9091}{mi}{99}\RG{7.628}{-0.003}{10.9091}{rm}{61}\RG{18.962}{-0.003}{10.9091}{rm}{48}%
}\hskip.1em\allowbreak{} 。

(3)假设\allowbreak{}\hskip.1em \RB{41.187}{8.686}{3.235}{%
\RG{0.000}{0.004}{10.9091}{mi}{114}\RG{5.111}{0.004}{10.9091}{rm}{40}\RG{9.265}{0.004}{10.9091}{mi}{67}\RG{17.690}{0.004}{10.9091}{rm}{41}\RG{24.868}{0.004}{10.9091}{rm}{61}\RG{36.214}{0.004}{10.9091}{rm}{49}%
}\hskip.1em\allowbreak{} ，仍写成\allowbreak{}\hskip.1em \RB{36.923}{8.529}{0.725}{%
\RG{0.000}{0.004}{10.9091}{mi}{67}\RG{11.457}{0.004}{10.9091}{rm}{61}\RG{22.795}{0.004}{10.9091}{mi}{117}\RG{28.871}{0.004}{10.9091}{mi}{118}\RG{34.430}{3.963}{7.9701}{sy}{48}%
}\hskip.1em\allowbreak{} ，其中\allowbreak{}\hskip.1em \RB{36.173}{8.315}{2.853}{%
\RG{0.000}{0.004}{10.9091}{mi}{117}\RG{6.076}{0.004}{10.9091}{mi}{59}\RG{10.855}{0.004}{10.9091}{mi}{118}\RG{19.451}{0.004}{10.9091}{sy}{54}\RG{19.451}{0.004}{10.9091}{rm}{61}\RG{30.786}{0.004}{10.9091}{rm}{48}%
}\hskip.1em\allowbreak{} 。题设给出

\[
\RB{247.517}{9.815}{3.242}{%
\RG{5.555}{4.500}{7.9701}{sy}{48}\RG{16.114}{4.500}{7.9701}{mi}{105}\RG{63.744}{4.500}{7.9701}{sy}{48}\RG{74.303}{4.500}{7.9701}{mi}{105}\RG{124.807}{4.500}{7.9701}{mi}{105}\RG{0.000}{-0.004}{10.9091}{mi}{118}\RG{8.106}{-0.004}{10.9091}{mi}{65}\RG{19.199}{-0.004}{10.9091}{mi}{117}\RG{28.309}{-0.004}{10.9091}{rm}{61}\RG{39.644}{-0.004}{10.9091}{rm}{116}\RG{43.800}{-0.004}{10.9091}{rm}{114}\RG{47.956}{-0.004}{10.9091}{rm}{40}\RG{52.110}{-0.004}{10.9091}{mi}{117}\RG{58.186}{-0.004}{10.9091}{mi}{118}\RG{66.295}{-0.004}{10.9091}{mi}{65}\RG{77.388}{-0.004}{10.9091}{rm}{41}\RG{84.567}{-0.004}{10.9091}{rm}{61}\RG{95.912}{-0.004}{10.9091}{rm}{116}\RG{100.069}{-0.004}{10.9091}{rm}{114}\RG{104.225}{-0.004}{10.9091}{rm}{40}\RG{108.374}{-0.004}{10.9091}{mi}{67}\RG{116.796}{-0.004}{10.9091}{mi}{65}\RG{127.892}{-0.004}{10.9091}{rm}{41}\RG{135.071}{-0.004}{10.9091}{rm}{61}\RG{146.416}{-0.004}{10.9091}{rm}{48}\RG{173.580}{-0.004}{10.9091}{rm}{40}\RG{177.737}{-0.004}{10.9091}{rm}{48}\RG{186.100}{-0.004}{10.9091}{sy}{20}\RG{197.616}{-0.004}{10.9091}{mi}{105}\RG{204.290}{-0.004}{10.9091}{sy}{20}\RG{215.805}{-0.004}{10.9091}{mi}{110}\RG{224.605}{-0.004}{10.9091}{sy}{0}\RG{235.515}{-0.004}{10.9091}{rm}{49}\RG{240.861}{-0.004}{10.9091}{rm}{41}\RG{245.010}{-0.004}{10.9091}{mi}{58}%
}
\]
设\allowbreak{}\hskip.1em \RB{96.238}{8.680}{3.240}{%
\RG{0.000}{-0.002}{10.9091}{mi}{102}\RG{6.430}{-0.002}{10.9091}{rm}{40}\RG{10.584}{-0.002}{10.9091}{mi}{116}\RG{14.441}{-0.002}{10.9091}{rm}{41}\RG{23.638}{-0.002}{10.9091}{rm}{61}\RG{37.002}{-0.002}{10.9091}{rm}{100}\RG{42.936}{-0.002}{10.9091}{rm}{101}\RG{47.682}{-0.002}{10.9091}{rm}{116}\RG{51.838}{-0.002}{10.9091}{rm}{40}\RG{55.994}{-0.002}{10.9091}{mi}{116}\RG{59.855}{-0.002}{10.9091}{mi}{73}\RG{64.562}{-1.638}{7.9701}{mi}{110}\RG{72.937}{-0.002}{10.9091}{sy}{0}\RG{84.655}{-0.002}{10.9091}{mi}{65}\RG{92.662}{-0.002}{10.9091}{rm}{41}%
}\hskip.1em\allowbreak{} 。先证明\allowbreak{}\hskip.1em \RB{78.163}{9.294}{2.607}{%
\RG{0.000}{-0.002}{10.9091}{mi}{117}\RG{6.076}{-0.002}{10.9091}{mi}{59}\RG{10.855}{-0.002}{10.9091}{mi}{65}\RG{18.862}{-0.002}{10.9091}{mi}{117}\RG{24.939}{-0.002}{10.9091}{mi}{59}\RG{29.728}{-0.002}{10.9091}{mi}{58}\RG{34.506}{-0.002}{10.9091}{mi}{58}\RG{39.295}{-0.002}{10.9091}{mi}{58}\RG{44.084}{-0.002}{10.9091}{mi}{59}\RG{48.863}{-0.002}{10.9091}{mi}{65}\RG{56.887}{3.958}{7.9701}{mi}{110}\RG{61.575}{3.958}{7.9701}{sy}{0}\RG{67.585}{3.958}{7.9701}{rm}{49}\RG{71.903}{-0.002}{10.9091}{mi}{117}%
}\hskip.1em\allowbreak{} 线性无关。若不然，存在非零多项式\allowbreak{}\hskip.1em \RB{37.790}{8.680}{3.229}{%
\RG{0.000}{-0.002}{10.9091}{mi}{103}\RG{8.537}{-0.002}{10.9091}{sy}{50}\RG{18.840}{-0.002}{10.9091}{mi}{75}\RG{28.710}{-0.002}{10.9091}{rm}{91}\RG{31.741}{-0.002}{10.9091}{mi}{116}\RG{35.599}{-0.002}{10.9091}{rm}{93}%
}\hskip.1em\allowbreak{} ，使得\allowbreak{}\hskip.1em \RB{44.417}{8.069}{2.738}{%
\RG{0.000}{-0.002}{10.9091}{rm}{100}\RG{5.935}{-0.002}{10.9091}{rm}{101}\RG{10.680}{-0.002}{10.9091}{rm}{103}\RG{17.986}{-0.002}{10.9091}{mi}{103}\RG{26.517}{-0.002}{10.9091}{mi}{60}\RG{37.862}{-0.002}{10.9091}{mi}{110}%
}\hskip.1em\allowbreak{} 且\allowbreak{}\hskip.1em \RB{47.654}{8.680}{3.240}{%
\RG{0.000}{-0.002}{10.9091}{mi}{103}\RG{5.506}{-0.002}{10.9091}{rm}{40}\RG{9.660}{-0.002}{10.9091}{mi}{65}\RG{17.668}{-0.002}{10.9091}{rm}{41}\RG{21.822}{-0.002}{10.9091}{mi}{117}\RG{30.931}{-0.002}{10.9091}{rm}{61}\RG{42.266}{-0.002}{10.9091}{rm}{48}%
}\hskip.1em\allowbreak{} 。由\allowbreak{}\hskip.1em \RB{6.456}{8.178}{2.727}{%
\RG{0.000}{-0.002}{10.9091}{mi}{102}%
}\hskip.1em\allowbreak{} 不可约，有\allowbreak{}\hskip.1em \RB{59.802}{8.680}{3.240}{%
\RG{0.000}{-0.002}{10.9091}{rm}{103}\RG{5.346}{-0.002}{10.9091}{rm}{99}\RG{10.091}{-0.002}{10.9091}{rm}{100}\RG{16.026}{-0.002}{10.9091}{rm}{40}\RG{20.177}{-0.002}{10.9091}{mi}{102}\RG{26.014}{-0.002}{10.9091}{mi}{59}\RG{30.803}{-0.002}{10.9091}{mi}{103}\RG{36.305}{-0.002}{10.9091}{rm}{41}\RG{43.484}{-0.002}{10.9091}{rm}{61}\RG{54.829}{-0.002}{10.9091}{rm}{49}%
}\hskip.1em\allowbreak{} ，于是多项式欧几里得算法给出\allowbreak{}\hskip.1em \RB{47.218}{8.679}{3.230}{%
\RG{0.000}{-0.003}{10.9091}{mi}{97}\RG{5.607}{-0.003}{10.9091}{mi}{59}\RG{10.386}{-0.003}{10.9091}{mi}{98}\RG{17.965}{-0.003}{10.9091}{sy}{50}\RG{28.268}{-0.003}{10.9091}{mi}{75}\RG{38.139}{-0.003}{10.9091}{rm}{91}\RG{41.170}{-0.003}{10.9091}{mi}{116}\RG{45.027}{-0.003}{10.9091}{rm}{93}%
}\hskip.1em\allowbreak{} ，满足\allowbreak{}\hskip.1em \RB{54.579}{8.177}{2.728}{%
\RG{0.000}{-0.003}{10.9091}{mi}{97}\RG{5.607}{-0.003}{10.9091}{mi}{102}\RG{14.461}{-0.003}{10.9091}{rm}{43}\RG{25.194}{-0.003}{10.9091}{mi}{98}\RG{29.732}{-0.003}{10.9091}{mi}{103}\RG{38.272}{-0.003}{10.9091}{rm}{61}\RG{49.606}{-0.003}{10.9091}{rm}{49}%
}\hskip.1em\allowbreak{} 。由\allowbreak{}\hskip.1em \RB{45.535}{8.679}{3.241}{%
\RG{0.000}{-0.003}{10.9091}{mi}{102}\RG{6.430}{-0.003}{10.9091}{rm}{40}\RG{10.584}{-0.003}{10.9091}{mi}{65}\RG{18.592}{-0.003}{10.9091}{rm}{41}\RG{25.770}{-0.003}{10.9091}{rm}{61}\RG{37.115}{-0.003}{10.9091}{mi}{79}%
}\hskip.1em\allowbreak{} 得

\[
\RB{155.449}{8.681}{3.239}{%
\RG{0.000}{-0.001}{10.9091}{mi}{117}\RG{9.110}{-0.001}{10.9091}{rm}{61}\RG{20.448}{-0.001}{10.9091}{mi}{97}\RG{26.056}{-0.001}{10.9091}{rm}{40}\RG{30.210}{-0.001}{10.9091}{mi}{65}\RG{38.218}{-0.001}{10.9091}{rm}{41}\RG{42.372}{-0.001}{10.9091}{mi}{102}\RG{48.801}{-0.001}{10.9091}{rm}{40}\RG{52.955}{-0.001}{10.9091}{mi}{65}\RG{60.963}{-0.001}{10.9091}{rm}{41}\RG{65.117}{-0.001}{10.9091}{mi}{117}\RG{73.621}{-0.001}{10.9091}{rm}{43}\RG{84.354}{-0.001}{10.9091}{mi}{98}\RG{88.895}{-0.001}{10.9091}{rm}{40}\RG{93.049}{-0.001}{10.9091}{mi}{65}\RG{101.056}{-0.001}{10.9091}{rm}{41}\RG{105.211}{-0.001}{10.9091}{mi}{103}\RG{110.717}{-0.001}{10.9091}{rm}{40}\RG{114.871}{-0.001}{10.9091}{mi}{65}\RG{122.879}{-0.001}{10.9091}{rm}{41}\RG{127.033}{-0.001}{10.9091}{mi}{117}\RG{136.142}{-0.001}{10.9091}{rm}{61}\RG{147.477}{-0.001}{10.9091}{rm}{48}\RG{152.822}{-0.001}{10.9091}{mi}{59}%
}
\]
与\allowbreak{}\hskip.1em \RB{25.833}{8.312}{2.856}{%
\RG{0.000}{0.001}{10.9091}{mi}{117}\RG{9.111}{0.001}{10.9091}{sy}{54}\RG{9.111}{0.001}{10.9091}{rm}{61}\RG{20.445}{0.001}{10.9091}{rm}{48}%
}\hskip.1em\allowbreak{} 矛盾。

因此\allowbreak{}\hskip.1em \RB{78.164}{9.293}{2.607}{%
\RG{0.000}{-0.002}{10.9091}{mi}{117}\RG{6.076}{-0.002}{10.9091}{mi}{59}\RG{10.855}{-0.002}{10.9091}{mi}{65}\RG{18.862}{-0.002}{10.9091}{mi}{117}\RG{24.939}{-0.002}{10.9091}{mi}{59}\RG{29.728}{-0.002}{10.9091}{mi}{58}\RG{34.506}{-0.002}{10.9091}{mi}{58}\RG{39.295}{-0.002}{10.9091}{mi}{58}\RG{44.084}{-0.002}{10.9091}{mi}{59}\RG{48.863}{-0.002}{10.9091}{mi}{65}\RG{56.887}{3.957}{7.9701}{mi}{110}\RG{61.575}{3.957}{7.9701}{sy}{0}\RG{67.585}{3.957}{7.9701}{rm}{49}\RG{71.904}{-0.002}{10.9091}{mi}{117}%
}\hskip.1em\allowbreak{} 是\allowbreak{}\hskip.1em \RB{14.823}{7.949}{0.502}{%
\RG{0.000}{-0.002}{10.9091}{mi}{75}\RG{9.872}{3.957}{7.9701}{mi}{110}%
}\hskip.1em\allowbreak{} 的一组基。上面的迹等式说明线性函数\allowbreak{}\hskip.1em \RB{37.253}{8.522}{0.622}{%
\RG{0.000}{-0.002}{10.9091}{mi}{120}\RG{9.100}{-0.002}{10.9091}{sy}{55}\RG{9.100}{-0.002}{10.9091}{sy}{33}\RG{23.040}{-0.002}{10.9091}{mi}{118}\RG{28.595}{3.957}{7.9701}{sy}{48}\RG{31.146}{-0.002}{10.9091}{mi}{120}%
}\hskip.1em\allowbreak{} 在这组基上全为零，故\allowbreak{}\hskip.1em \RB{27.858}{8.522}{0.732}{%
\RG{0.000}{-0.003}{10.9091}{mi}{118}\RG{5.556}{3.957}{7.9701}{sy}{48}\RG{11.136}{-0.003}{10.9091}{rm}{61}\RG{22.471}{-0.003}{10.9091}{rm}{48}%
}\hskip.1em\allowbreak{} ，与\allowbreak{}\hskip.1em \RB{25.308}{8.308}{2.859}{%
\RG{0.000}{-0.003}{10.9091}{mi}{118}\RG{8.586}{-0.003}{10.9091}{sy}{54}\RG{8.586}{-0.003}{10.9091}{rm}{61}\RG{19.921}{-0.003}{10.9091}{rm}{48}%
}\hskip.1em\allowbreak{} 矛盾。因此\allowbreak{}\hskip.1em \RB{41.183}{8.679}{3.241}{%
\RG{0.000}{-0.003}{10.9091}{mi}{114}\RG{5.111}{-0.003}{10.9091}{rm}{40}\RG{9.265}{-0.003}{10.9091}{mi}{67}\RG{17.691}{-0.003}{10.9091}{rm}{41}\RG{24.875}{-0.003}{10.9091}{sy}{54}\RG{24.875}{-0.003}{10.9091}{rm}{61}\RG{36.210}{-0.003}{10.9091}{rm}{49}%
}\hskip.1em\allowbreak{} 。

小结：不可约的特征多项式使每个非零向量都成为循环向量；秩一矩阵的迹条件可转化为一个非零线性函数在循环基上恒为零。

% Original page 3.
\section*{2018S06}
请用高代教材[1]第六章的方法证明以下问题：(1)设数域\allowbreak{}\hskip.1em \RB{10.035}{7.950}{0.501}{%
\RG{0.000}{-0.001}{10.9091}{mi}{75}%
}\hskip.1em\allowbreak{} 上的\allowbreak{}\hskip.1em \RB{6.555}{5.310}{0.610}{%
\RG{0.000}{-0.001}{10.9091}{mi}{110}%
}\hskip.1em\allowbreak{} 阶方阵\allowbreak{}\hskip.1em \RB{8.202}{8.288}{0.501}{%
\RG{0.000}{-0.001}{10.9091}{mi}{65}%
}\hskip.1em\allowbreak{} 满足\allowbreak{}\hskip.1em \RB{39.503}{8.288}{2.210}{%
\RG{0.000}{-0.001}{10.9091}{mi}{65}\RG{8.007}{3.958}{7.9701}{mi}{109}\RG{18.415}{-0.001}{10.9091}{rm}{61}\RG{29.841}{-0.001}{10.9091}{mi}{73}\RG{34.552}{-1.637}{7.9701}{mi}{110}%
}\hskip.1em\allowbreak{} ，其中\allowbreak{}\hskip.1em \RB{9.522}{5.310}{0.610}{%
\RG{0.000}{-0.001}{10.9091}{mi}{109}%
}\hskip.1em\allowbreak{} 是正整数。证明：\allowbreak{}\hskip.1em \RB{8.202}{8.288}{0.501}{%
\RG{0.000}{-0.001}{10.9091}{mi}{65}%
}\hskip.1em\allowbreak{} 在复数域上

可对角化；(2)设\allowbreak{}\hskip.1em \RB{8.715}{7.947}{0.733}{%
\RG{0.000}{-0.004}{10.9091}{mi}{86}%
}\hskip.1em\allowbreak{} 是数域\allowbreak{}\hskip.1em \RB{10.035}{7.947}{0.504}{%
\RG{0.000}{-0.004}{10.9091}{mi}{75}%
}\hskip.1em\allowbreak{} 上的\allowbreak{}\hskip.1em \RB{6.555}{5.307}{0.613}{%
\RG{0.000}{-0.004}{10.9091}{mi}{110}%
}\hskip.1em\allowbreak{} 维线性空间，\allowbreak{}\hskip.1em \RB{7.122}{5.307}{2.860}{%
\RG{0.000}{-0.004}{10.9091}{mi}{39}%
}\hskip.1em\allowbreak{} 是\allowbreak{}\hskip.1em \RB{8.715}{7.947}{0.733}{%
\RG{0.000}{-0.004}{10.9091}{mi}{86}%
}\hskip.1em\allowbreak{} 上的线性变换，满足\allowbreak{}\hskip.1em \RB{39.780}{7.947}{2.860}{%
\RG{0.000}{-0.004}{10.9091}{mi}{39}\RG{6.997}{3.955}{7.9701}{mi}{109}\RG{17.318}{-0.004}{10.9091}{rm}{61}\RG{28.656}{-0.004}{10.9091}{mi}{73}\RG{33.367}{-1.640}{7.9701}{mi}{86}%
}\hskip.1em\allowbreak{} ，其中\allowbreak{}\hskip.1em \RB{9.522}{5.307}{0.613}{%
\RG{0.000}{-0.004}{10.9091}{mi}{109}%
}\hskip.1em\allowbreak{}

是正整数。设

\[
\RB{206.919}{17.151}{11.525}{%
\RG{175.802}{11.814}{7.9701}{mi}{109}\RG{182.637}{11.814}{7.9701}{sy}{0}\RG{188.647}{11.814}{7.9701}{rm}{49}\RG{177.590}{9.088}{9.0910}{ex}{88}\RG{165.551}{7.377}{10.9091}{rm}{49}\RG{201.327}{4.501}{7.9701}{mi}{105}\RG{0.000}{-0.004}{10.9091}{mi}{87}\RG{14.598}{-0.004}{10.9091}{rm}{61}\RG{25.937}{-0.004}{10.9091}{sy}{102}\RG{31.391}{-0.004}{10.9091}{mi}{118}\RG{39.977}{-0.004}{10.9091}{sy}{50}\RG{50.280}{-0.004}{10.9091}{mi}{86}\RG{61.912}{-0.004}{10.9091}{rm}{58}\RG{67.909}{-0.004}{10.9091}{mi}{39}\RG{74.907}{-0.004}{10.9091}{rm}{40}\RG{79.061}{-0.004}{10.9091}{mi}{118}\RG{84.616}{-0.004}{10.9091}{rm}{41}\RG{91.795}{-0.004}{10.9091}{rm}{61}\RG{103.139}{-0.004}{10.9091}{mi}{118}\RG{108.695}{-0.004}{10.9091}{sy}{103}\RG{114.150}{-0.004}{10.9091}{mi}{59}\RG{140.747}{-0.004}{10.9091}{mi}{32}\RG{151.119}{-0.004}{10.9091}{rm}{61}\RG{194.330}{-0.004}{10.9091}{mi}{39}\RG{204.412}{-0.004}{10.9091}{mi}{58}\RG{163.550}{-7.487}{10.9091}{mi}{109}\RG{177.904}{-10.873}{7.9701}{mi}{105}\RG{180.535}{-10.873}{7.9701}{rm}{61}\RG{186.542}{-10.873}{7.9701}{rm}{48}\RR{163.550}{2.724}{9.343}{0.364}%
}
\]
证明：\allowbreak{}\hskip.1em \RB{69.689}{8.679}{3.241}{%
\RG{0.000}{-0.003}{10.9091}{rm}{116}\RG{4.156}{-0.003}{10.9091}{rm}{114}\RG{8.313}{-0.003}{10.9091}{rm}{40}\RG{12.462}{-0.003}{10.9091}{mi}{32}\RG{19.798}{-0.003}{10.9091}{rm}{41}\RG{26.976}{-0.003}{10.9091}{rm}{61}\RG{38.322}{-0.003}{10.9091}{rm}{100}\RG{44.256}{-0.003}{10.9091}{rm}{105}\RG{47.224}{-0.003}{10.9091}{rm}{109}\RG{57.996}{-0.003}{10.9091}{mi}{87}%
}\hskip.1em\allowbreak{} 。

证明：(1)令\allowbreak{}\hskip.1em \RB{53.193}{9.909}{3.623}{%
\RG{0.000}{-0.004}{10.9091}{mi}{16}\RG{4.666}{-1.640}{7.9701}{mi}{106}\RG{11.695}{-0.004}{10.9091}{rm}{61}\RG{23.034}{-0.004}{10.9091}{mi}{101}\RG{27.984}{3.955}{7.9701}{rm}{50}\RG{31.848}{3.955}{7.9701}{mi}{25}\RG{36.546}{3.955}{7.9701}{rm}{105}\RG{38.692}{3.955}{7.9701}{mi}{106}\RG{42.234}{3.955}{7.9701}{mi}{61}\RG{46.096}{3.955}{7.9701}{mi}{109}%
}\hskip.1em\allowbreak{} ，\allowbreak{}\hskip.1em \RB{48.806}{7.751}{2.729}{%
\RG{0.000}{-0.004}{10.9091}{rm}{48}\RG{8.371}{-0.004}{10.9091}{sy}{20}\RG{19.887}{-0.004}{10.9091}{mi}{106}\RG{27.938}{-0.004}{10.9091}{mi}{60}\RG{39.284}{-0.004}{10.9091}{mi}{109}%
}\hskip.1em\allowbreak{} ，并定义

\[
\RB{86.360}{17.151}{11.768}{%
\RG{38.596}{11.814}{7.9701}{mi}{109}\RG{45.431}{11.814}{7.9701}{sy}{0}\RG{51.441}{11.814}{7.9701}{rm}{49}\RG{40.385}{9.087}{9.0910}{ex}{88}\RG{28.344}{7.376}{10.9091}{rm}{49}\RG{62.613}{4.500}{7.9701}{sy}{0}\RG{68.623}{4.500}{7.9701}{mi}{96}\RG{80.242}{4.500}{7.9701}{mi}{96}\RG{0.000}{-0.005}{10.9091}{mi}{80}\RG{13.914}{-0.005}{10.9091}{rm}{61}\RG{57.123}{-0.005}{10.9091}{mi}{16}\RG{72.234}{-0.005}{10.9091}{mi}{65}\RG{83.853}{-0.005}{10.9091}{mi}{58}\RG{6.885}{-1.641}{7.9701}{mi}{106}\RG{61.789}{-2.701}{7.9701}{mi}{106}\RG{26.343}{-7.487}{10.9091}{mi}{109}\RG{40.435}{-11.115}{7.9701}{mi}{96}\RG{43.592}{-11.115}{7.9701}{rm}{61}\RG{49.599}{-11.115}{7.9701}{rm}{48}\RR{26.343}{2.723}{9.343}{0.364}%
}
\]
由\allowbreak{}\hskip.1em \RB{39.329}{8.293}{2.205}{%
\RG{0.000}{0.004}{10.9091}{mi}{65}\RG{8.008}{3.963}{7.9701}{mi}{109}\RG{18.328}{0.004}{10.9091}{rm}{61}\RG{29.667}{0.004}{10.9091}{mi}{73}\RG{34.378}{-1.632}{7.9701}{mi}{110}%
}\hskip.1em\allowbreak{} 及\allowbreak{}\hskip.1em \RB{32.117}{8.097}{4.676}{%
\RG{0.000}{0.004}{10.9091}{mi}{16}\RG{4.666}{-2.693}{7.9701}{mi}{106}\RG{5.490}{3.963}{7.9701}{mi}{109}\RG{15.809}{0.004}{10.9091}{rm}{61}\RG{27.144}{0.004}{10.9091}{rm}{49}%
}\hskip.1em\allowbreak{} ，有

\[
\RB{124.707}{15.522}{11.703}{%
\RG{49.752}{11.814}{7.9701}{mi}{109}\RG{46.604}{9.087}{9.0910}{ex}{88}\RG{34.691}{7.377}{10.9091}{mi}{16}\RG{39.357}{5.740}{7.9701}{mi}{106}\RG{67.043}{4.501}{7.9701}{sy}{0}\RG{73.053}{4.501}{7.9701}{mi}{96}\RG{84.672}{4.501}{7.9701}{mi}{96}\RG{0.000}{-0.004}{10.9091}{mi}{65}\RG{8.007}{-0.004}{10.9091}{mi}{80}\RG{21.921}{-0.004}{10.9091}{rm}{61}\RG{61.554}{-0.004}{10.9091}{mi}{16}\RG{76.664}{-0.004}{10.9091}{mi}{65}\RG{91.313}{-0.004}{10.9091}{rm}{61}\RG{102.652}{-0.004}{10.9091}{mi}{16}\RG{111.317}{-0.004}{10.9091}{mi}{80}\RG{122.200}{-0.004}{10.9091}{mi}{58}\RG{14.892}{-1.641}{7.9701}{mi}{106}\RG{107.317}{-1.641}{7.9701}{mi}{106}\RG{118.201}{-1.641}{7.9701}{mi}{106}\RG{66.219}{-2.701}{7.9701}{mi}{106}\RG{34.352}{-7.487}{10.9091}{mi}{109}\RG{46.654}{-11.115}{7.9701}{mi}{96}\RG{49.811}{-11.115}{7.9701}{rm}{61}\RG{55.819}{-11.115}{7.9701}{rm}{49}\RR{34.352}{2.724}{9.343}{0.364}%
}
\]
有限等比数列求和给出

\noindent \RB{245.772}{29.033}{4.676}{%
\RG{48.007}{28.542}{9.0910}{ex}{56}\RG{0.000}{21.270}{7.9701}{mi}{109}\RG{6.836}{21.270}{7.9701}{sy}{0}\RG{12.845}{21.270}{7.9701}{rm}{49}\RG{189.630}{21.270}{7.9701}{mi}{109}\RG{196.465}{21.270}{7.9701}{sy}{0}\RG{202.476}{21.270}{7.9701}{rm}{49}\RG{48.007}{20.361}{9.0910}{ex}{60}\RG{1.789}{18.542}{9.0910}{ex}{88}\RG{191.419}{18.542}{9.0910}{ex}{88}\RG{56.089}{17.399}{10.9091}{mi}{109}\RG{65.438}{17.399}{10.9091}{mi}{59}\RG{79.303}{17.399}{10.9091}{mi}{96}\RG{86.821}{17.399}{10.9091}{rm}{61}\RG{98.156}{17.399}{10.9091}{rm}{48}\RG{103.501}{17.399}{10.9091}{mi}{59}\RG{24.017}{13.956}{7.9701}{sy}{0}\RG{30.027}{13.956}{7.9701}{mi}{96}\RG{18.528}{9.451}{10.9091}{mi}{16}\RG{36.668}{9.451}{10.9091}{rm}{61}\RT{155.084}{9.451}{10.9091}{故}\RG{208.157}{9.451}{10.9091}{mi}{80}\RG{222.072}{9.451}{10.9091}{rm}{61}\RG{233.410}{9.451}{10.9091}{mi}{73}\RG{243.265}{9.451}{10.9091}{mi}{58}\RG{215.043}{7.815}{7.9701}{mi}{106}\RG{238.121}{7.815}{7.9701}{mi}{110}\RG{23.193}{6.754}{7.9701}{mi}{106}\RG{48.007}{3.997}{9.0910}{ex}{58}\RG{56.089}{-0.002}{10.9091}{rm}{48}\RG{61.429}{-0.002}{10.9091}{mi}{59}\RG{79.308}{-0.002}{10.9091}{rm}{49}\RG{87.680}{-0.002}{10.9091}{sy}{20}\RG{99.195}{-0.002}{10.9091}{mi}{96}\RG{106.700}{-0.002}{10.9091}{mi}{60}\RG{118.046}{-0.002}{10.9091}{mi}{109}\RG{127.395}{-0.002}{10.9091}{mi}{59}\RG{1.646}{-1.417}{7.9701}{mi}{106}\RG{5.190}{-1.417}{7.9701}{rm}{61}\RG{11.198}{-1.417}{7.9701}{rm}{48}\RG{191.275}{-1.417}{7.9701}{mi}{106}\RG{194.820}{-1.417}{7.9701}{rm}{61}\RG{200.827}{-1.417}{7.9701}{rm}{48}%
}

于是每个\allowbreak{}\hskip.1em \RB{32.233}{8.184}{0.933}{%
\RG{0.000}{0.004}{10.9091}{mi}{120}\RG{9.100}{0.004}{10.9091}{sy}{50}\RG{19.404}{0.004}{10.9091}{bb}{67}\RG{27.282}{3.962}{7.9701}{mi}{110}%
}\hskip.1em\allowbreak{} 都是\allowbreak{}\hskip.1em \RB{71.357}{7.955}{2.602}{%
\RG{0.000}{0.004}{10.9091}{mi}{80}\RG{6.885}{-1.633}{7.9701}{rm}{48}\RG{11.203}{0.004}{10.9091}{mi}{120}\RG{17.268}{0.004}{10.9091}{mi}{59}\RG{22.047}{0.004}{10.9091}{mi}{58}\RG{26.836}{0.004}{10.9091}{mi}{58}\RG{31.625}{0.004}{10.9091}{mi}{58}\RG{36.403}{0.004}{10.9091}{mi}{59}\RG{41.193}{0.004}{10.9091}{mi}{80}\RG{48.085}{-1.633}{7.9701}{mi}{109}\RG{54.920}{-1.633}{7.9701}{sy}{0}\RG{60.930}{-1.633}{7.9701}{rm}{49}\RG{65.249}{0.004}{10.9091}{mi}{120}%
}\hskip.1em\allowbreak{} 之和，且\allowbreak{}\hskip.1em \RB{120.357}{8.686}{3.617}{%
\RG{0.000}{0.004}{10.9091}{mi}{80}\RG{6.885}{-1.633}{7.9701}{mi}{106}\RG{10.884}{0.004}{10.9091}{mi}{120}\RG{19.984}{0.004}{10.9091}{sy}{50}\RG{30.287}{0.004}{10.9091}{mi}{69}\RG{38.194}{-1.633}{7.9701}{mi}{106}\RG{45.223}{0.004}{10.9091}{rm}{58}\RG{48.191}{0.004}{10.9091}{rm}{61}\RG{59.525}{0.004}{10.9091}{rm}{107}\RG{64.882}{0.004}{10.9091}{rm}{101}\RG{69.627}{0.004}{10.9091}{rm}{114}\RG{73.784}{0.004}{10.9091}{rm}{40}\RG{77.926}{0.004}{10.9091}{mi}{65}\RG{87.855}{0.004}{10.9091}{sy}{0}\RG{98.262}{0.004}{10.9091}{mi}{16}\RG{102.928}{-1.633}{7.9701}{mi}{106}\RG{106.926}{0.004}{10.9091}{mi}{73}\RG{111.638}{-1.633}{7.9701}{mi}{110}\RG{116.781}{0.004}{10.9091}{rm}{41}%
}\hskip.1em\allowbreak{} 。若\allowbreak{}\hskip.1em \RB{31.217}{7.933}{2.722}{%
\RG{0.000}{0.004}{10.9091}{mi}{121}\RG{8.630}{0.004}{10.9091}{sy}{50}\RG{18.933}{0.004}{10.9091}{mi}{69}\RG{26.841}{-1.633}{7.9701}{mi}{107}%
}\hskip.1em\allowbreak{} ，再次由等比数列求和得

\[
\RB{169.534}{29.040}{4.670}{%
\RG{114.035}{28.549}{9.0910}{ex}{56}\RG{62.723}{25.458}{9.0910}{ex}{32}\RG{81.440}{25.458}{9.0910}{ex}{33}\RG{88.637}{22.942}{7.9701}{mi}{96}\RG{44.196}{21.276}{7.9701}{mi}{109}\RG{51.031}{21.276}{7.9701}{sy}{0}\RG{57.041}{21.276}{7.9701}{rm}{49}\RG{114.035}{20.367}{9.0910}{ex}{60}\RG{45.985}{18.549}{9.0910}{ex}{88}\RG{122.116}{17.405}{10.9091}{mi}{121}\RG{127.702}{17.405}{10.9091}{mi}{59}\RG{141.578}{17.405}{10.9091}{mi}{106}\RG{149.655}{17.405}{10.9091}{rm}{61}\RG{160.994}{17.405}{10.9091}{mi}{107}\RG{166.907}{17.405}{10.9091}{mi}{59}\RG{33.945}{16.838}{10.9091}{rm}{49}\RG{71.012}{16.838}{10.9091}{mi}{16}\RG{75.677}{15.202}{7.9701}{mi}{107}\RG{0.000}{9.457}{10.9091}{mi}{80}\RG{10.884}{9.457}{10.9091}{mi}{121}\RG{19.514}{9.457}{10.9091}{rm}{61}\RG{94.066}{9.457}{10.9091}{mi}{121}\RG{102.696}{9.457}{10.9091}{rm}{61}\RG{6.885}{7.821}{7.9701}{mi}{106}\RG{114.035}{4.003}{9.0910}{ex}{58}\RG{31.943}{1.975}{10.9091}{mi}{109}\RG{71.348}{1.975}{10.9091}{mi}{16}\RG{76.013}{0.338}{7.9701}{mi}{106}\RG{122.116}{0.005}{10.9091}{rm}{48}\RG{127.457}{0.005}{10.9091}{mi}{59}\RG{141.585}{0.005}{10.9091}{mi}{106}\RG{149.655}{0.005}{10.9091}{sy}{54}\RG{149.655}{0.005}{10.9091}{rm}{61}\RG{160.994}{0.005}{10.9091}{mi}{107}\RG{166.907}{0.005}{10.9091}{mi}{58}\RG{46.035}{-1.653}{7.9701}{mi}{96}\RG{49.192}{-1.653}{7.9701}{rm}{61}\RG{55.199}{-1.653}{7.9701}{rm}{48}\RR{31.943}{12.185}{9.343}{0.364}\RR{71.012}{12.185}{9.337}{0.364}%
}
\]
因此，对\allowbreak{}\hskip.1em \RB{45.917}{7.772}{2.847}{%
\RG{0.000}{7.272}{9.0910}{ex}{80}\RG{9.596}{-2.274}{7.9701}{mi}{107}\RG{16.087}{-0.001}{10.9091}{mi}{121}\RG{21.294}{-1.637}{7.9701}{mi}{107}\RG{29.096}{-0.001}{10.9091}{rm}{61}\RG{40.529}{-0.001}{10.9091}{rm}{48}%
}\hskip.1em\allowbreak{} （\allowbreak{}\hskip.1em \RB{35.697}{7.928}{2.726}{%
\RG{0.000}{-0.001}{10.9091}{mi}{121}\RG{5.209}{-1.637}{7.9701}{mi}{107}\RG{13.010}{-0.001}{10.9091}{sy}{50}\RG{23.413}{-0.001}{10.9091}{mi}{69}\RG{31.320}{-1.637}{7.9701}{mi}{107}%
}\hskip.1em\allowbreak{} ）作用\allowbreak{}\hskip.1em \RB{10.389}{7.950}{3.621}{%
\RG{0.000}{-0.001}{10.9091}{mi}{80}\RG{6.885}{-1.637}{7.9701}{mi}{106}%
}\hskip.1em\allowbreak{} 即得\allowbreak{}\hskip.1em \RB{29.157}{7.754}{3.621}{%
\RG{0.000}{-0.001}{10.9091}{mi}{121}\RG{5.209}{-1.637}{7.9701}{mi}{106}\RG{12.337}{-0.001}{10.9091}{rm}{61}\RG{23.770}{-0.001}{10.9091}{rm}{48}%
}\hskip.1em\allowbreak{} ，所以\allowbreak{}\hskip.1em \RB{68.084}{9.729}{4.681}{%
\RG{0.000}{-0.001}{10.9091}{bb}{67}\RG{7.879}{3.958}{7.9701}{mi}{110}\RG{16.151}{-0.001}{10.9091}{rm}{61}\RG{27.590}{7.272}{9.0910}{ex}{76}\RG{37.691}{4.393}{7.9701}{mi}{109}\RG{37.691}{-2.697}{7.9701}{mi}{106}\RG{41.236}{-2.697}{7.9701}{rm}{61}\RG{44.526}{4.393}{7.9701}{sy}{0}\RG{47.243}{-2.697}{7.9701}{rm}{48}\RG{50.537}{4.393}{7.9701}{rm}{49}\RG{56.673}{-0.001}{10.9091}{mi}{69}\RG{64.580}{-1.637}{7.9701}{mi}{106}%
}\hskip.1em\allowbreak{} 。各非零\allowbreak{}\hskip.1em \RB{11.411}{7.928}{3.621}{%
\RG{0.000}{-0.001}{10.9091}{mi}{69}\RG{7.907}{-1.637}{7.9701}{mi}{106}%
}\hskip.1em\allowbreak{} 的基合在一起构成特征向量基，从而\allowbreak{}\hskip.1em \RB{8.202}{8.288}{0.501}{%
\RG{0.000}{-0.001}{10.9091}{mi}{65}%
}\hskip.1em\allowbreak{} 可对角化。

(2)由\allowbreak{}\hskip.1em \RB{39.780}{7.947}{2.860}{%
\RG{0.000}{-0.004}{10.9091}{mi}{39}\RG{6.997}{3.955}{7.9701}{mi}{109}\RG{17.318}{-0.004}{10.9091}{rm}{61}\RG{28.657}{-0.004}{10.9091}{mi}{73}\RG{33.367}{-1.641}{7.9701}{mi}{86}%
}\hskip.1em\allowbreak{} ，有

\[
\RB{91.406}{15.526}{11.441}{%
\RG{45.194}{11.819}{7.9701}{mi}{109}\RG{42.046}{9.091}{9.0910}{ex}{88}\RG{31.795}{7.380}{10.9091}{rm}{49}\RG{63.994}{4.505}{7.9701}{mi}{105}\RG{0.000}{-0.000}{10.9091}{mi}{39}\RG{6.993}{-0.000}{10.9091}{mi}{32}\RG{17.364}{-0.000}{10.9091}{rm}{61}\RG{56.996}{0.001}{10.9091}{mi}{39}\RG{70.109}{0.000}{10.9091}{rm}{61}\RG{81.447}{0.001}{10.9091}{mi}{32}\RG{88.778}{0.001}{10.9091}{mi}{59}\RG{29.794}{-7.483}{10.9091}{mi}{109}\RG{42.360}{-10.868}{7.9701}{mi}{105}\RG{44.991}{-10.868}{7.9701}{rm}{61}\RG{50.998}{-10.868}{7.9701}{rm}{49}\RR{29.794}{2.727}{9.343}{0.364}%
}
\]
故\allowbreak{}\hskip.1em \RB{48.201}{8.067}{2.729}{%
\RG{0.000}{-0.004}{10.9091}{rm}{73}\RG{3.851}{-0.004}{10.9091}{rm}{109}\RG{14.626}{-0.004}{10.9091}{mi}{32}\RG{24.992}{-0.004}{10.9091}{sy}{18}\RG{36.508}{-0.004}{10.9091}{mi}{87}%
}\hskip.1em\allowbreak{} 。对每个\allowbreak{}\hskip.1em \RB{32.955}{7.947}{0.940}{%
\RG{0.000}{-0.004}{10.9091}{mi}{119}\RG{10.959}{-0.004}{10.9091}{sy}{50}\RG{21.262}{-0.004}{10.9091}{mi}{87}%
}\hskip.1em\allowbreak{} ，有\allowbreak{}\hskip.1em \RB{48.553}{9.269}{3.242}{%
\RG{0.000}{-0.004}{10.9091}{mi}{39}\RG{6.997}{3.955}{7.9701}{mi}{105}\RG{10.082}{-0.004}{10.9091}{rm}{40}\RG{14.236}{-0.004}{10.9091}{mi}{119}\RG{22.166}{-0.004}{10.9091}{rm}{41}\RG{29.344}{-0.004}{10.9091}{rm}{61}\RG{40.689}{-0.004}{10.9091}{mi}{119}%
}\hskip.1em\allowbreak{} ，因而\allowbreak{}\hskip.1em \RB{45.806}{8.678}{3.242}{%
\RG{0.000}{-0.004}{10.9091}{mi}{32}\RG{7.337}{-0.004}{10.9091}{rm}{40}\RG{11.490}{-0.004}{10.9091}{mi}{119}\RG{19.419}{-0.004}{10.9091}{rm}{41}\RG{26.597}{-0.004}{10.9091}{rm}{61}\RG{37.942}{-0.004}{10.9091}{mi}{119}%
}\hskip.1em\allowbreak{} ，所以\allowbreak{}\hskip.1em \RB{48.024}{8.067}{2.729}{%
\RG{0.000}{-0.004}{10.9091}{rm}{73}\RG{3.851}{-0.004}{10.9091}{rm}{109}\RG{14.626}{-0.004}{10.9091}{mi}{32}\RG{24.992}{-0.004}{10.9091}{rm}{61}\RG{36.331}{-0.004}{10.9091}{mi}{87}%
}\hskip.1em\allowbreak{} 且\allowbreak{}\hskip.1em \RB{33.287}{9.291}{2.729}{%
\RG{0.000}{-0.004}{10.9091}{mi}{32}\RG{7.336}{3.955}{7.9701}{rm}{50}\RG{14.685}{-0.004}{10.9091}{rm}{61}\RG{26.023}{-0.004}{10.9091}{mi}{32}%
}\hskip.1em\allowbreak{} 。

任意\allowbreak{}\hskip.1em \RB{32.957}{7.954}{0.934}{%
\RG{0.000}{0.003}{10.9091}{mi}{118}\RG{11.262}{0.003}{10.9091}{sy}{50}\RG{24.242}{0.003}{10.9091}{mi}{86}%
}\hskip.1em\allowbreak{} 可写为\allowbreak{}\hskip.1em \RB{111.737}{8.685}{3.235}{%
\RG{0.000}{0.003}{10.9091}{mi}{118}\RG{11.262}{0.003}{10.9091}{rm}{61}\RG{25.278}{0.003}{10.9091}{mi}{32}\RG{32.613}{0.003}{10.9091}{rm}{40}\RG{36.767}{0.003}{10.9091}{mi}{118}\RG{42.323}{0.003}{10.9091}{rm}{41}\RG{49.971}{0.003}{10.9091}{rm}{43}\RG{61.774}{0.003}{10.9091}{rm}{40}\RG{65.929}{0.003}{10.9091}{mi}{118}\RG{74.979}{0.003}{10.9091}{sy}{0}\RG{86.959}{0.003}{10.9091}{mi}{32}\RG{94.294}{0.003}{10.9091}{rm}{40}\RG{98.448}{0.003}{10.9091}{mi}{118}\RG{104.005}{0.003}{10.9091}{rm}{41}\RG{108.161}{0.003}{10.9091}{rm}{41}%
}\hskip.1em\allowbreak{} ，其中两项分别属于\allowbreak{}\hskip.1em \RB{11.693}{7.954}{0.726}{%
\RG{0.000}{0.003}{10.9091}{mi}{87}%
}\hskip.1em\allowbreak{} 与\allowbreak{}\hskip.1em \RB{23.324}{8.074}{2.723}{%
\RG{0.000}{0.003}{10.9091}{rm}{107}\RG{5.346}{0.003}{10.9091}{rm}{101}\RG{10.091}{0.003}{10.9091}{rm}{114}\RG{16.060}{0.003}{10.9091}{mi}{32}%
}\hskip.1em\allowbreak{} 。若\allowbreak{}\hskip.1em \RB{68.276}{8.073}{2.724}{%
\RG{0.000}{0.002}{10.9091}{mi}{119}\RG{10.959}{0.002}{10.9091}{sy}{50}\RG{21.262}{0.002}{10.9091}{mi}{87}\RG{35.255}{0.002}{10.9091}{sy}{92}\RG{44.952}{0.002}{10.9091}{rm}{107}\RG{50.298}{0.002}{10.9091}{rm}{101}\RG{55.043}{0.002}{10.9091}{rm}{114}\RG{61.012}{0.002}{10.9091}{mi}{32}%
}\hskip.1em\allowbreak{} ，则\allowbreak{}\hskip.1em \RB{65.629}{8.684}{3.236}{%
\RG{0.000}{0.002}{10.9091}{mi}{119}\RG{10.960}{0.002}{10.9091}{rm}{61}\RG{22.298}{0.002}{10.9091}{mi}{32}\RG{29.635}{0.002}{10.9091}{rm}{40}\RG{33.789}{0.002}{10.9091}{mi}{119}\RG{41.717}{0.002}{10.9091}{rm}{41}\RG{48.896}{0.002}{10.9091}{rm}{61}\RG{60.241}{0.002}{10.9091}{rm}{48}%
}\hskip.1em\allowbreak{} ，故

\[
\RB{73.770}{8.075}{2.722}{%
\RG{0.000}{0.004}{10.9091}{mi}{86}\RG{11.631}{0.004}{10.9091}{rm}{61}\RG{22.970}{0.004}{10.9091}{mi}{87}\RG{36.962}{0.004}{10.9091}{sy}{8}\RG{47.871}{0.004}{10.9091}{rm}{107}\RG{53.217}{0.004}{10.9091}{rm}{101}\RG{57.962}{0.004}{10.9091}{rm}{114}\RG{63.932}{0.004}{10.9091}{mi}{32}\RG{71.263}{0.004}{10.9091}{mi}{58}%
}
\]
令\allowbreak{}\hskip.1em \RB{51.287}{8.067}{0.733}{%
\RG{0.000}{-0.004}{10.9091}{mi}{100}\RG{8.581}{-0.004}{10.9091}{rm}{61}\RG{19.916}{-0.004}{10.9091}{rm}{100}\RG{25.851}{-0.004}{10.9091}{rm}{105}\RG{28.818}{-0.004}{10.9091}{rm}{109}\RG{39.594}{-0.004}{10.9091}{mi}{87}%
}\hskip.1em\allowbreak{} ，在与此直和分解相适应的基下，\allowbreak{}\hskip.1em \RB{7.264}{8.067}{2.729}{%
\RG{0.000}{-0.004}{10.9091}{mi}{32}%
}\hskip.1em\allowbreak{} 的表示矩阵为\allowbreak{}\hskip.1em \RB{49.845}{8.678}{3.242}{%
\RG{0.000}{-0.004}{10.9091}{rm}{100}\RG{5.935}{-0.004}{10.9091}{rm}{105}\RG{8.902}{-0.004}{10.9091}{rm}{97}\RG{14.248}{-0.004}{10.9091}{rm}{103}\RG{19.735}{-0.004}{10.9091}{rm}{40}\RG{23.882}{-0.004}{10.9091}{mi}{73}\RG{28.594}{-1.640}{7.9701}{mi}{100}\RG{33.024}{-0.004}{10.9091}{mi}{59}\RG{37.802}{-0.004}{10.9091}{mi}{79}\RG{46.269}{-0.004}{10.9091}{rm}{41}%
}\hskip.1em\allowbreak{} ，因此\allowbreak{}\hskip.1em \RB{35.189}{8.678}{3.242}{%
\RG{0.000}{-0.004}{10.9091}{rm}{116}\RG{4.156}{-0.004}{10.9091}{rm}{114}\RG{8.313}{-0.004}{10.9091}{rm}{40}\RG{12.462}{-0.004}{10.9091}{mi}{32}\RG{19.798}{-0.004}{10.9091}{rm}{41}\RG{26.976}{-0.004}{10.9091}{rm}{61}%
}\hskip.1em\allowbreak{}

\[
\RB{62.571}{10.097}{1.813}{%
\RG{0.000}{-0.004}{10.9091}{mi}{100}\RG{8.580}{-0.004}{10.9091}{rm}{61}\RG{19.915}{-0.004}{10.9091}{rm}{100}\RG{25.850}{-0.004}{10.9091}{rm}{105}\RG{28.817}{-0.004}{10.9091}{rm}{109}\RG{39.594}{-0.004}{10.9091}{mi}{87}\RT{51.162}{-0.004}{10.9091}{。}%
}
\]
小结：单位根加权平均给出各特征子空间上的投影；不加权的平均算子则是到不动点空间的投影。

% Original page 4.
\section*{2018S07}
设\allowbreak{}\hskip.1em \RB{8.202}{8.290}{0.500}{%
\RG{0.000}{0.001}{10.9091}{mi}{65}%
}\hskip.1em\allowbreak{} 为\allowbreak{}\hskip.1em \RB{6.555}{5.312}{0.608}{%
\RG{0.000}{0.001}{10.9091}{mi}{110}%
}\hskip.1em\allowbreak{} 阶复矩阵，请用高代教材[1]第六章的方法证明：

\noindent \RB{242.664}{10.099}{18.936}{%
\RT{0.000}{-0.001}{10.9091}{(}\RT{4.156}{-0.001}{10.9091}{1}\RT{9.502}{-0.001}{10.9091}{)}\RG{19.113}{-0.001}{10.9091}{mi}{65}\RT{30.678}{-0.001}{10.9091}{可}\RT{41.587}{-0.001}{10.9091}{对}\RT{52.496}{-0.001}{10.9091}{角}\RT{63.406}{-0.001}{10.9091}{化}\RT{74.315}{-0.001}{10.9091}{当}\RT{85.224}{-0.001}{10.9091}{且}\RT{96.133}{-0.001}{10.9091}{仅}\RT{107.043}{-0.001}{10.9091}{当}\RG{121.507}{-0.001}{10.9091}{mi}{65}\RT{133.072}{-0.001}{10.9091}{的}\RT{143.981}{-0.001}{10.9091}{极}\RT{154.890}{-0.001}{10.9091}{小}\RT{165.800}{-0.001}{10.9091}{多}\RT{176.709}{-0.001}{10.9091}{项}\RT{187.618}{-0.001}{10.9091}{式}\RT{198.527}{-0.001}{10.9091}{无}\RT{209.437}{-0.001}{10.9091}{重}\RT{220.346}{-0.001}{10.9091}{根}\RT{231.255}{-0.001}{10.9091}{；}\RG{86.533}{-9.925}{7.9701}{sy}{48}\RG{111.460}{-9.925}{7.9701}{sy}{48}\RT{0.000}{-16.319}{10.9091}{(}\RT{4.156}{-16.319}{10.9091}{2}\RT{9.502}{-16.319}{10.9091}{)}\RT{19.113}{-16.319}{10.9091}{若}\RG{33.579}{-16.319}{10.9091}{mi}{65}\RT{45.143}{-16.319}{10.9091}{满}\RT{56.052}{-16.319}{10.9091}{足}\RG{70.518}{-16.319}{10.9091}{mi}{65}\RG{78.525}{-16.319}{10.9091}{mi}{65}\RG{92.114}{-16.319}{10.9091}{rm}{61}\RG{103.452}{-16.319}{10.9091}{mi}{65}\RG{114.010}{-16.319}{10.9091}{mi}{65}\RT{122.018}{-16.319}{10.9091}{，}\RT{132.928}{-16.319}{10.9091}{则}\RG{147.394}{-16.319}{10.9091}{mi}{65}\RT{158.957}{-16.319}{10.9091}{可}\RT{169.866}{-16.319}{10.9091}{对}\RT{180.776}{-16.319}{10.9091}{角}\RT{191.685}{-16.319}{10.9091}{化}\RT{202.594}{-16.319}{10.9091}{。}\RR{78.525}{-7.592}{8.007}{0.364}\RR{103.452}{-7.592}{8.007}{0.364}%
}

证明：(1)设\allowbreak{}\hskip.1em \RB{27.174}{8.678}{3.242}{%
\RG{0.000}{-0.004}{10.9091}{mi}{109}\RG{9.344}{-1.641}{7.9701}{mi}{65}\RG{15.587}{-0.004}{10.9091}{rm}{40}\RG{19.740}{-0.004}{10.9091}{mi}{116}\RG{23.598}{-0.004}{10.9091}{rm}{41}%
}\hskip.1em\allowbreak{} 为\allowbreak{}\hskip.1em \RB{8.202}{8.286}{0.504}{%
\RG{0.000}{-0.004}{10.9091}{mi}{65}%
}\hskip.1em\allowbreak{} 的极小多项式。若\allowbreak{}\hskip.1em \RB{8.202}{8.286}{0.504}{%
\RG{0.000}{-0.004}{10.9091}{mi}{65}%
}\hskip.1em\allowbreak{} 可对角化，其互异特征值为\allowbreak{}\hskip.1em \RB{44.400}{8.067}{2.609}{%
\RG{0.000}{-0.004}{10.9091}{mi}{21}\RG{6.232}{-1.641}{7.9701}{rm}{49}\RG{10.550}{-0.004}{10.9091}{mi}{59}\RG{15.328}{-0.004}{10.9091}{mi}{58}\RG{20.117}{-0.004}{10.9091}{mi}{58}\RG{24.895}{-0.004}{10.9091}{mi}{58}\RG{29.684}{-0.004}{10.9091}{mi}{59}\RG{34.474}{-0.004}{10.9091}{mi}{21}\RG{40.707}{-1.641}{7.9701}{mi}{115}%
}\hskip.1em\allowbreak{} ，则对任意\allowbreak{}\hskip.1em \RB{36.439}{8.687}{3.223}{%
\RG{0.000}{0.005}{10.9091}{mi}{104}\RG{9.179}{0.005}{10.9091}{sy}{50}\RG{19.482}{0.005}{10.9091}{bb}{67}\RG{27.361}{0.005}{10.9091}{rm}{91}\RG{30.391}{0.005}{10.9091}{mi}{116}\RG{34.248}{0.005}{10.9091}{rm}{93}%
}\hskip.1em\allowbreak{} ，有

\[
\RB{203.791}{8.679}{3.623}{%
\RG{0.000}{-0.003}{10.9091}{mi}{104}\RG{6.149}{-0.003}{10.9091}{rm}{40}\RG{10.303}{-0.003}{10.9091}{mi}{65}\RG{18.310}{-0.003}{10.9091}{rm}{41}\RG{25.489}{-0.003}{10.9091}{rm}{61}\RG{36.834}{-0.003}{10.9091}{mi}{79}\RG{59.233}{-0.003}{10.9091}{sy}{40}\RG{68.331}{-0.003}{10.9091}{sy}{41}\RG{93.173}{-0.003}{10.9091}{mi}{104}\RG{99.322}{-0.003}{10.9091}{rm}{40}\RG{103.476}{-0.003}{10.9091}{mi}{21}\RG{113.706}{-0.003}{10.9091}{rm}{41}\RG{120.885}{-0.003}{10.9091}{rm}{61}\RG{132.230}{-0.003}{10.9091}{rm}{48}\RG{148.485}{-0.003}{10.9091}{rm}{40}\RG{152.642}{-0.003}{10.9091}{rm}{49}\RG{161.005}{-0.003}{10.9091}{sy}{20}\RG{172.520}{-0.003}{10.9091}{mi}{106}\RG{180.583}{-0.003}{10.9091}{sy}{20}\RG{192.099}{-0.003}{10.9091}{mi}{115}\RG{197.130}{-0.003}{10.9091}{rm}{41}\RG{201.284}{-0.003}{10.9091}{mi}{58}\RG{109.707}{-1.639}{7.9701}{mi}{106}%
}
\]
\noindent \RB{285.346}{26.415}{4.685}{%
\RG{67.495}{23.588}{9.0910}{ex}{81}\RG{76.081}{20.709}{7.9701}{mi}{115}\RT{0.000}{16.315}{10.9091}{因}\RT{10.909}{16.315}{10.9091}{此}\RG{25.375}{16.315}{10.9091}{mi}{109}\RG{40.961}{16.315}{10.9091}{rm}{40}\RG{45.115}{16.315}{10.9091}{mi}{116}\RG{48.972}{16.315}{10.9091}{rm}{41}\RG{56.150}{16.315}{10.9091}{rm}{61}\RG{89.954}{16.315}{10.9091}{rm}{40}\RG{94.107}{16.315}{10.9091}{mi}{116}\RG{100.389}{16.315}{10.9091}{sy}{0}\RG{111.298}{16.315}{10.9091}{mi}{21}\RG{121.528}{16.315}{10.9091}{rm}{41}\RT{125.682}{16.315}{10.9091}{，}\RT{136.591}{16.315}{10.9091}{它}\RT{147.501}{16.315}{10.9091}{没}\RT{158.410}{16.315}{10.9091}{有}\RT{169.319}{16.315}{10.9091}{重}\RT{180.229}{16.315}{10.9091}{根}\RT{191.138}{16.315}{10.9091}{。}\RG{34.718}{14.679}{7.9701}{mi}{65}\RG{117.530}{14.679}{7.9701}{mi}{106}\RG{76.081}{13.618}{7.9701}{mi}{106}\RG{79.625}{13.618}{7.9701}{rm}{61}\RG{85.632}{13.618}{7.9701}{rm}{49}\RG{111.132}{7.269}{9.0910}{ex}{81}\RG{119.718}{4.390}{7.9701}{mi}{115}\RT{21.818}{-0.004}{10.9091}{反}\RT{32.727}{-0.004}{10.9091}{之}\RT{43.637}{-0.004}{10.9091}{，}\RT{54.546}{-0.004}{10.9091}{设}\RG{69.012}{-0.004}{10.9091}{mi}{109}\RG{84.598}{-0.004}{10.9091}{rm}{40}\RG{88.752}{-0.004}{10.9091}{mi}{116}\RG{92.609}{-0.004}{10.9091}{rm}{41}\RG{99.787}{-0.004}{10.9091}{rm}{61}\RG{133.591}{-0.004}{10.9091}{rm}{40}\RG{137.745}{-0.004}{10.9091}{mi}{116}\RG{144.026}{-0.004}{10.9091}{sy}{0}\RG{154.936}{-0.004}{10.9091}{mi}{21}\RG{165.165}{-0.004}{10.9091}{rm}{41}\RT{169.319}{-0.004}{10.9091}{，}\RT{180.229}{-0.004}{10.9091}{其}\RT{191.138}{-0.004}{10.9091}{中}\RG{205.604}{-0.004}{10.9091}{mi}{21}\RT{219.390}{-0.004}{10.9091}{两}\RT{230.299}{-0.004}{10.9091}{两}\RT{241.209}{-0.004}{10.9091}{不}\RT{252.118}{-0.004}{10.9091}{同}\RT{263.027}{-0.004}{10.9091}{。}\RT{273.936}{-0.004}{10.9091}{令}\RG{78.355}{-1.640}{7.9701}{mi}{65}\RG{161.166}{-1.640}{7.9701}{mi}{106}\RG{211.834}{-1.640}{7.9701}{mi}{106}\RG{119.718}{-2.701}{7.9701}{mi}{106}\RG{123.262}{-2.701}{7.9701}{rm}{61}\RG{129.270}{-2.701}{7.9701}{rm}{49}%
}

\[
\RB{89.760}{15.449}{13.178}{%
\RG{36.093}{9.089}{9.0910}{ex}{89}\RG{54.882}{7.378}{10.9091}{mi}{116}\RG{61.164}{7.378}{10.9091}{sy}{0}\RG{72.073}{7.378}{10.9091}{mi}{21}\RG{78.304}{5.741}{7.9701}{mi}{107}\RG{0.000}{-0.003}{10.9091}{mi}{96}\RG{8.481}{-0.003}{10.9091}{rm}{40}\RG{12.635}{-0.003}{10.9091}{mi}{116}\RG{16.492}{-0.003}{10.9091}{rm}{41}\RG{23.670}{-0.003}{10.9091}{rm}{61}\RG{87.253}{-0.003}{10.9091}{mi}{58}\RG{4.482}{-1.639}{7.9701}{mi}{106}\RG{51.696}{-7.486}{10.9091}{mi}{21}\RG{64.349}{-7.486}{10.9091}{sy}{0}\RG{75.259}{-7.486}{10.9091}{mi}{21}\RG{57.927}{-9.122}{7.9701}{mi}{106}\RG{81.490}{-9.122}{7.9701}{mi}{107}\RG{35.015}{-11.114}{7.9701}{mi}{107}\RG{39.233}{-11.114}{7.9701}{sy}{54}\RG{39.233}{-11.114}{7.9701}{rm}{61}\RG{45.242}{-11.114}{7.9701}{mi}{106}\RR{51.696}{2.725}{34.467}{0.364}%
}
\]
多项式\allowbreak{}\hskip.1em \RB{54.365}{8.680}{4.258}{%
\RG{0.000}{7.272}{9.0910}{ex}{80}\RG{9.596}{-2.274}{7.9701}{mi}{106}\RG{15.413}{-0.002}{10.9091}{mi}{96}\RG{19.896}{-1.638}{7.9701}{mi}{106}\RG{23.894}{-0.002}{10.9091}{rm}{40}\RG{28.048}{-0.002}{10.9091}{mi}{116}\RG{31.905}{-0.002}{10.9091}{rm}{41}\RG{38.483}{-0.002}{10.9091}{sy}{0}\RG{49.392}{-0.002}{10.9091}{rm}{49}%
}\hskip.1em\allowbreak{} 的次数小于\allowbreak{}\hskip.1em \RB{5.006}{5.309}{0.611}{%
\RG{0.000}{-0.002}{10.9091}{mi}{115}%
}\hskip.1em\allowbreak{} ，且在\allowbreak{}\hskip.1em \RB{44.400}{8.069}{2.607}{%
\RG{0.000}{-0.002}{10.9091}{mi}{21}\RG{6.231}{-1.638}{7.9701}{rm}{49}\RG{10.550}{-0.002}{10.9091}{mi}{59}\RG{15.328}{-0.002}{10.9091}{mi}{58}\RG{20.117}{-0.002}{10.9091}{mi}{58}\RG{24.895}{-0.002}{10.9091}{mi}{58}\RG{29.684}{-0.002}{10.9091}{mi}{59}\RG{34.474}{-0.002}{10.9091}{mi}{21}\RG{40.707}{-1.638}{7.9701}{mi}{115}%
}\hskip.1em\allowbreak{} 处全为零，故它是零多项式。因此

\[
\RB{272.858}{16.057}{12.857}{%
\RG{4.922}{11.813}{7.9701}{mi}{115}\RG{0.143}{9.086}{9.0910}{ex}{88}\RG{197.692}{7.376}{10.9091}{mi}{109}\RG{213.278}{7.375}{10.9091}{rm}{40}\RG{217.432}{7.375}{10.9091}{mi}{65}\RG{225.440}{7.375}{10.9091}{rm}{41}\RG{207.035}{5.739}{7.9701}{mi}{65}\RG{15.236}{-0.005}{10.9091}{mi}{96}\RG{23.717}{-0.005}{10.9091}{rm}{40}\RG{27.871}{-0.005}{10.9091}{mi}{65}\RG{35.878}{-0.005}{10.9091}{rm}{41}\RG{43.056}{-0.005}{10.9091}{rm}{61}\RG{54.401}{-0.005}{10.9091}{mi}{73}\RG{64.256}{-0.005}{10.9091}{mi}{59}\RG{90.860}{-0.005}{10.9091}{rm}{40}\RG{95.014}{-0.005}{10.9091}{mi}{65}\RG{105.445}{-0.005}{10.9091}{sy}{0}\RG{116.355}{-0.005}{10.9091}{mi}{21}\RG{126.585}{-0.005}{10.9091}{mi}{73}\RG{136.439}{-0.005}{10.9091}{rm}{41}\RG{140.593}{-0.005}{10.9091}{mi}{96}\RG{149.075}{-0.005}{10.9091}{rm}{40}\RG{153.228}{-0.005}{10.9091}{mi}{65}\RG{161.236}{-0.005}{10.9091}{rm}{41}\RG{168.414}{-0.005}{10.9091}{rm}{61}\RG{250.558}{-0.005}{10.9091}{rm}{61}\RG{261.896}{-0.005}{10.9091}{mi}{79}\RG{270.350}{-0.005}{10.9091}{mi}{58}\RG{180.850}{-0.215}{9.0910}{ex}{81}\RG{19.718}{-1.642}{7.9701}{mi}{106}\RG{59.113}{-1.642}{7.9701}{mi}{110}\RG{122.586}{-1.642}{7.9701}{mi}{106}\RG{131.296}{-1.642}{7.9701}{mi}{110}\RG{145.076}{-1.642}{7.9701}{mi}{106}\RG{203.662}{-7.488}{10.9091}{rm}{40}\RG{207.816}{-7.488}{10.9091}{mi}{21}\RG{220.470}{-7.488}{10.9091}{sy}{0}\RG{231.379}{-7.488}{10.9091}{mi}{21}\RG{242.282}{-7.488}{10.9091}{rm}{41}\RG{214.047}{-9.124}{7.9701}{mi}{106}\RG{237.610}{-9.124}{7.9701}{mi}{107}\RG{189.436}{-9.761}{7.9701}{mi}{107}\RG{193.653}{-9.761}{7.9701}{sy}{54}\RG{193.653}{-9.761}{7.9701}{rm}{61}\RG{199.663}{-9.761}{7.9701}{mi}{106}\RG{0.000}{-10.873}{7.9701}{mi}{106}\RG{3.544}{-10.873}{7.9701}{rm}{61}\RG{9.552}{-10.873}{7.9701}{rm}{49}\RR{180.850}{2.723}{65.586}{0.364}%
}
\]
每个向量均为各特征子空间中向量之和。若\allowbreak{}\hskip.1em \RB{46.544}{7.775}{4.254}{%
\RG{0.000}{7.275}{9.0910}{ex}{80}\RG{9.596}{-2.270}{7.9701}{mi}{106}\RG{15.412}{0.002}{10.9091}{mi}{120}\RG{21.482}{-1.634}{7.9701}{mi}{106}\RG{29.168}{0.002}{10.9091}{rm}{61}\RG{41.157}{0.002}{10.9091}{rm}{48}%
}\hskip.1em\allowbreak{} ，其中\allowbreak{}\hskip.1em \RB{53.560}{8.291}{3.618}{%
\RG{0.000}{0.002}{10.9091}{mi}{65}\RG{8.007}{0.002}{10.9091}{mi}{120}\RG{14.078}{-1.634}{7.9701}{mi}{106}\RG{21.763}{0.002}{10.9091}{rm}{61}\RG{33.757}{0.002}{10.9091}{mi}{21}\RG{39.988}{-1.634}{7.9701}{mi}{106}\RG{43.987}{0.002}{10.9091}{mi}{120}\RG{50.057}{-1.634}{7.9701}{mi}{106}%
}\hskip.1em\allowbreak{} ，则作用\allowbreak{}\hskip.1em \RB{24.893}{8.684}{3.236}{%
\RG{0.000}{0.002}{10.9091}{mi}{96}\RG{4.482}{-1.634}{7.9701}{mi}{107}\RG{9.154}{0.002}{10.9091}{rm}{40}\RG{13.308}{0.002}{10.9091}{mi}{65}\RG{21.316}{0.002}{10.9091}{rm}{41}%
}\hskip.1em\allowbreak{}后，由\allowbreak{}\hskip.1em \RB{87.350}{8.684}{3.619}{%
\RG{0.000}{0.002}{10.9091}{mi}{96}\RG{4.483}{-1.635}{7.9701}{mi}{107}\RG{9.154}{0.002}{10.9091}{rm}{40}\RG{13.309}{0.002}{10.9091}{mi}{65}\RG{21.316}{0.002}{10.9091}{rm}{41}\RG{25.470}{0.002}{10.9091}{mi}{120}\RG{31.540}{-1.635}{7.9701}{mi}{106}\RG{38.658}{0.002}{10.9091}{rm}{61}\RG{50.084}{0.002}{10.9091}{mi}{96}\RG{54.566}{-1.635}{7.9701}{mi}{107}\RG{59.239}{0.002}{10.9091}{rm}{40}\RG{63.392}{0.002}{10.9091}{mi}{21}\RG{69.623}{-1.635}{7.9701}{mi}{106}\RG{73.622}{0.002}{10.9091}{rm}{41}\RG{77.776}{0.002}{10.9091}{mi}{120}\RG{83.846}{-1.635}{7.9701}{mi}{106}%
}\hskip.1em\allowbreak{} 得\allowbreak{}\hskip.1em \RB{30.670}{7.756}{2.208}{%
\RG{0.000}{0.002}{10.9091}{mi}{120}\RG{6.070}{-1.635}{7.9701}{mi}{107}\RG{13.860}{0.002}{10.9091}{rm}{61}\RG{25.282}{0.002}{10.9091}{rm}{48}%
}\hskip.1em\allowbreak{} 。故这些特征子空间的和为直和，从中选取基即可得到特征向量基，因而\allowbreak{}\hskip.1em \RB{8.202}{8.290}{0.500}{%
\RG{0.000}{0.001}{10.9091}{mi}{65}%
}\hskip.1em\allowbreak{} 可对角化。

(2)对任意\allowbreak{}\hskip.1em \RB{27.537}{8.178}{0.938}{%
\RG{0.000}{-0.002}{10.9091}{mi}{21}\RG{9.261}{-0.002}{10.9091}{sy}{50}\RG{19.564}{-0.002}{10.9091}{bb}{67}%
}\hskip.1em\allowbreak{} ，令\allowbreak{}\hskip.1em \RB{59.279}{8.287}{2.211}{%
\RG{0.000}{-0.002}{10.9091}{mi}{84}\RG{10.775}{-0.002}{10.9091}{rm}{61}\RG{22.113}{-0.002}{10.9091}{mi}{65}\RG{32.511}{-0.002}{10.9091}{sy}{0}\RG{43.386}{-0.002}{10.9091}{mi}{21}\RG{49.616}{-0.002}{10.9091}{mi}{73}\RG{54.328}{-1.638}{7.9701}{mi}{110}%
}\hskip.1em\allowbreak{} 。由题设得\allowbreak{}\hskip.1em \RB{50.753}{10.958}{0.502}{%
\RG{0.000}{-0.002}{10.9091}{mi}{84}\RG{7.744}{-0.002}{10.9091}{mi}{84}\RG{15.489}{6.392}{7.9701}{sy}{48}\RG{21.070}{-0.002}{10.9091}{rm}{61}\RG{32.409}{-0.002}{10.9091}{mi}{84}\RG{40.153}{6.392}{7.9701}{sy}{48}\RG{42.704}{-0.002}{10.9091}{mi}{84}\RR{7.744}{8.725}{7.744}{0.364}\RR{32.409}{8.725}{7.744}{0.364}%
}\hskip.1em\allowbreak{} 。记\allowbreak{}\hskip.1em \RB{100.238}{9.294}{4.258}{%
\RG{0.000}{-0.002}{10.9091}{sy}{107}\RG{5.455}{-0.002}{10.9091}{mi}{120}\RG{11.525}{-0.002}{10.9091}{sy}{107}\RG{16.979}{3.958}{7.9701}{rm}{50}\RG{24.328}{-0.002}{10.9091}{rm}{61}\RG{35.666}{-0.002}{10.9091}{mi}{120}\RG{41.736}{3.958}{7.9701}{sy}{48}\RG{44.287}{-0.002}{10.9091}{mi}{120}\RG{53.387}{-0.002}{10.9091}{rm}{61}\RG{64.726}{7.271}{9.0910}{ex}{80}\RG{74.322}{-2.274}{7.9701}{mi}{106}\RG{80.140}{-0.002}{10.9091}{sy}{106}\RG{83.170}{-0.002}{10.9091}{mi}{120}\RG{89.240}{-1.638}{7.9701}{mi}{106}\RG{93.239}{-0.002}{10.9091}{sy}{106}\RG{96.269}{3.958}{7.9701}{rm}{50}\RR{35.666}{5.968}{6.070}{0.364}%
}\hskip.1em\allowbreak{} ，则

\[
\RB{173.527}{10.959}{3.227}{%
\RG{61.594}{6.394}{7.9701}{sy}{48}\RG{116.438}{6.394}{7.9701}{sy}{48}\RG{152.626}{6.394}{7.9701}{sy}{48}\RG{24.724}{4.504}{7.9701}{rm}{50}\RG{49.481}{4.504}{7.9701}{sy}{48}\RG{98.398}{4.504}{7.9701}{sy}{48}\RG{166.701}{4.504}{7.9701}{rm}{50}\RG{43.411}{0.000}{10.9091}{mi}{120}\RG{53.850}{0.000}{10.9091}{mi}{84}\RG{64.145}{0.000}{10.9091}{mi}{84}\RG{71.880}{0.000}{10.9091}{mi}{120}\RG{80.989}{0.000}{10.9091}{rm}{61}\RG{92.328}{0.000}{10.9091}{mi}{120}\RG{100.949}{0.000}{10.9091}{mi}{84}\RG{108.693}{0.000}{10.9091}{mi}{84}\RG{118.988}{0.000}{10.9091}{mi}{120}\RG{128.088}{0.000}{10.9091}{rm}{61}\RG{139.426}{0.000}{10.9091}{sy}{107}\RG{144.881}{0.000}{10.9091}{mi}{84}\RG{155.176}{0.000}{10.9091}{mi}{120}\RG{161.246}{0.000}{10.9091}{sy}{107}\RG{171.019}{0.000}{10.9091}{mi}{58}\RG{0.000}{-0.000}{10.9091}{sy}{107}\RG{5.455}{-0.000}{10.9091}{mi}{84}\RG{13.190}{-0.000}{10.9091}{mi}{120}\RG{19.270}{-0.000}{10.9091}{sy}{107}\RG{32.072}{-0.000}{10.9091}{rm}{61}\RR{43.411}{5.970}{6.070}{0.364}\RR{53.850}{8.727}{7.744}{0.364}\RR{92.328}{5.970}{6.070}{0.364}\RR{108.693}{8.727}{7.744}{0.364}\RR{144.881}{8.727}{7.744}{0.364}%
}
\]
所以\allowbreak{}\hskip.1em \RB{33.566}{7.900}{0.725}{%
\RG{0.000}{0.004}{10.9091}{mi}{84}\RG{7.735}{0.004}{10.9091}{mi}{120}\RG{16.844}{0.004}{10.9091}{rm}{61}\RG{28.179}{0.004}{10.9091}{rm}{48}%
}\hskip.1em\allowbreak{} 蕴含\allowbreak{}\hskip.1em \RB{36.118}{10.963}{0.725}{%
\RG{0.000}{0.004}{10.9091}{mi}{84}\RG{7.744}{6.398}{7.9701}{sy}{48}\RG{10.295}{0.004}{10.9091}{mi}{120}\RG{19.395}{0.004}{10.9091}{rm}{61}\RG{30.730}{0.004}{10.9091}{rm}{48}\RR{0.000}{8.731}{7.744}{0.364}%
}\hskip.1em\allowbreak{} ，对后式取\allowbreak{}\hskip.1em \RB{11.561}{0.500}{23.178}{%
\RG{0.000}{-19.951}{10.9091}{sy}{107}\RG{5.454}{-19.951}{10.9091}{mi}{120}%
}\hskip.1em\allowbreak{} 共\allowbreak{}\hskip.1em \RB{9.423}{0.500}{23.178}{%
\RG{0.000}{-19.951}{10.9091}{sy}{107}\RG{5.454}{-15.447}{7.9701}{rm}{50}%
}\hskip.1em\allowbreak{} 轭\allowbreak{}\hskip.1em \RB{8.213}{0.500}{18.946}{%
\RG{0.000}{-19.951}{10.9091}{rm}{61}%
}\hskip.1em\allowbreak{} 转\allowbreak{}\hskip.1em \RB{6.107}{0.500}{20.560}{%
\RG{0.000}{-19.951}{10.9091}{mi}{120}\RR{0.000}{-13.982}{6.070}{0.364}%
}\hskip.1em\allowbreak{} 置\allowbreak{}\hskip.1em \RB{10.600}{0.500}{20.451}{%
\RG{0.000}{-15.447}{7.9701}{sy}{48}\RG{2.550}{-19.951}{10.9091}{mi}{84}%
}\hskip.1em\allowbreak{} 得\allowbreak{}\hskip.1em \RB{48.688}{8.528}{22.677}{%
\RG{0.000}{-19.951}{10.9091}{mi}{121}\RG{8.629}{-19.951}{10.9091}{rm}{61}\RG{12.571}{0.004}{10.9091}{mi}{120}\RG{18.641}{3.962}{7.9701}{sy}{48}\RG{19.963}{-19.951}{10.9091}{rm}{48}\RG{21.191}{0.004}{10.9091}{mi}{84}\RG{25.308}{-19.951}{10.9091}{mi}{59}\RG{31.966}{0.004}{10.9091}{rm}{61}\RG{43.301}{0.004}{10.9091}{rm}{48}\RR{12.571}{5.973}{6.070}{0.364}%
}\hskip.1em\allowbreak{} 。若同时\allowbreak{}\hskip.1em \RB{33.900}{7.900}{2.722}{%
\RG{0.000}{0.004}{10.9091}{mi}{120}\RG{9.100}{0.004}{10.9091}{rm}{61}\RG{20.439}{0.004}{10.9091}{mi}{84}\RG{28.185}{0.004}{10.9091}{mi}{121}%
}\hskip.1em\allowbreak{} ，则从而\allowbreak{}\hskip.1em \RB{25.822}{7.753}{0.731}{%
\RG{0.000}{-0.002}{10.9091}{mi}{120}\RG{9.100}{-0.002}{10.9091}{rm}{61}\RG{20.434}{-0.002}{10.9091}{rm}{48}%
}\hskip.1em\allowbreak{} 。因此

\[
\RB{258.547}{9.291}{3.229}{%
\RG{0.000}{-0.002}{10.9091}{rm}{107}\RG{5.346}{-0.002}{10.9091}{rm}{101}\RG{10.091}{-0.002}{10.9091}{rm}{114}\RG{16.060}{-0.002}{10.9091}{mi}{84}\RG{26.229}{-0.002}{10.9091}{sy}{92}\RG{35.926}{-0.002}{10.9091}{rm}{73}\RG{39.777}{-0.002}{10.9091}{rm}{109}\RG{50.552}{-0.002}{10.9091}{mi}{84}\RG{61.326}{-0.002}{10.9091}{rm}{61}\RG{72.665}{-0.002}{10.9091}{sy}{102}\RG{78.120}{-0.002}{10.9091}{rm}{48}\RG{83.461}{-0.002}{10.9091}{sy}{103}\RG{88.915}{-0.002}{10.9091}{mi}{58}\RT{244.388}{-0.002}{10.9091}{(}\RT{248.545}{-0.002}{10.9091}{1}\RT{253.890}{-0.002}{10.9091}{)}%
}
\]
若\allowbreak{}\hskip.1em \RB{15.400}{5.308}{2.140}{%
\RG{0.000}{-0.003}{10.9091}{mi}{109}\RG{9.343}{-1.640}{7.9701}{mi}{65}%
}\hskip.1em\allowbreak{} 有重根\allowbreak{}\hskip.1em \RB{6.336}{8.068}{0.623}{%
\RG{0.000}{-0.003}{10.9091}{mi}{21}%
}\hskip.1em\allowbreak{} ，写成\allowbreak{}\hskip.1em \RB{95.748}{9.292}{3.241}{%
\RG{0.000}{-0.003}{10.9091}{mi}{109}\RG{9.344}{-1.640}{7.9701}{mi}{65}\RG{15.587}{-0.003}{10.9091}{rm}{40}\RG{19.740}{-0.003}{10.9091}{mi}{116}\RG{23.598}{-0.003}{10.9091}{rm}{41}\RG{31.092}{-0.003}{10.9091}{rm}{61}\RG{42.743}{-0.003}{10.9091}{rm}{40}\RG{46.892}{-0.003}{10.9091}{mi}{116}\RG{53.297}{-0.003}{10.9091}{sy}{0}\RG{64.329}{-0.003}{10.9091}{mi}{21}\RG{70.561}{-0.003}{10.9091}{rm}{41}\RG{74.715}{3.956}{7.9701}{rm}{50}\RG{79.033}{-0.003}{10.9091}{mi}{113}\RG{84.160}{-0.003}{10.9091}{rm}{40}\RG{88.315}{-0.003}{10.9091}{mi}{116}\RG{92.172}{-0.003}{10.9091}{rm}{41}%
}\hskip.1em\allowbreak{} 。由于\allowbreak{}\hskip.1em \RB{48.692}{8.679}{3.241}{%
\RG{0.000}{-0.003}{10.9091}{rm}{40}\RG{4.154}{-0.003}{10.9091}{mi}{116}\RG{10.559}{-0.003}{10.9091}{sy}{0}\RG{21.592}{-0.003}{10.9091}{mi}{21}\RG{27.823}{-0.003}{10.9091}{rm}{41}\RG{31.977}{-0.003}{10.9091}{mi}{113}\RG{37.105}{-0.003}{10.9091}{rm}{40}\RG{41.258}{-0.003}{10.9091}{mi}{116}\RG{45.116}{-0.003}{10.9091}{rm}{41}%
}\hskip.1em\allowbreak{} 的次数小于\allowbreak{}\hskip.1em \RB{33.385}{8.068}{2.739}{%
\RG{0.000}{-0.003}{10.9091}{rm}{100}\RG{5.935}{-0.003}{10.9091}{rm}{101}\RG{10.680}{-0.003}{10.9091}{rm}{103}\RG{17.985}{-0.003}{10.9091}{mi}{109}\RG{27.329}{-1.640}{7.9701}{mi}{65}%
}\hskip.1em\allowbreak{} ，它不能零化\allowbreak{}\hskip.1em \RB{8.202}{8.286}{0.504}{%
\RG{0.000}{-0.004}{10.9091}{mi}{65}%
}\hskip.1em\allowbreak{} ，故存在\allowbreak{}\hskip.1em \RB{5.715}{5.307}{2.729}{%
\RG{0.000}{-0.004}{10.9091}{mi}{121}%
}\hskip.1em\allowbreak{} ，使

\[
\RB{115.434}{8.679}{3.241}{%
\RG{0.000}{-0.003}{10.9091}{mi}{120}\RG{9.100}{-0.003}{10.9091}{rm}{61}\RG{20.435}{-0.003}{10.9091}{rm}{40}\RG{24.593}{-0.003}{10.9091}{mi}{65}\RG{35.025}{-0.003}{10.9091}{sy}{0}\RG{45.934}{-0.003}{10.9091}{mi}{21}\RG{52.163}{-0.003}{10.9091}{mi}{73}\RG{62.020}{-0.003}{10.9091}{rm}{41}\RG{66.174}{-0.003}{10.9091}{mi}{113}\RG{71.302}{-0.003}{10.9091}{rm}{40}\RG{75.456}{-0.003}{10.9091}{mi}{65}\RG{83.463}{-0.003}{10.9091}{rm}{41}\RG{87.617}{-0.003}{10.9091}{mi}{121}\RG{96.247}{-0.003}{10.9091}{sy}{54}\RG{96.247}{-0.003}{10.9091}{rm}{61}\RG{107.582}{-0.003}{10.9091}{rm}{48}\RG{112.927}{-0.003}{10.9091}{mi}{58}\RG{56.877}{-1.640}{7.9701}{mi}{110}%
}
\]
此时\allowbreak{}\hskip.1em \RB{42.558}{7.956}{0.932}{%
\RG{0.000}{0.005}{10.9091}{mi}{120}\RG{9.340}{0.005}{10.9091}{sy}{50}\RG{19.882}{0.005}{10.9091}{rm}{73}\RG{23.733}{0.005}{10.9091}{rm}{109}\RG{34.508}{0.005}{10.9091}{mi}{84}%
}\hskip.1em\allowbreak{} ，而\allowbreak{}\hskip.1em \RB{86.395}{8.687}{3.234}{%
\RG{0.000}{0.005}{10.9091}{mi}{84}\RG{7.735}{0.005}{10.9091}{mi}{120}\RG{17.084}{0.005}{10.9091}{rm}{61}\RG{28.662}{0.005}{10.9091}{mi}{109}\RG{38.005}{-1.632}{7.9701}{mi}{65}\RG{44.248}{0.005}{10.9091}{rm}{40}\RG{48.402}{0.005}{10.9091}{mi}{65}\RG{56.410}{0.005}{10.9091}{rm}{41}\RG{60.563}{0.005}{10.9091}{mi}{121}\RG{69.433}{0.005}{10.9091}{rm}{61}\RG{81.008}{0.005}{10.9091}{rm}{48}%
}\hskip.1em\allowbreak{} ，与(1)矛盾。故\allowbreak{}\hskip.1em \RB{15.400}{5.316}{2.132}{%
\RG{0.000}{0.005}{10.9091}{mi}{109}\RG{9.343}{-1.632}{7.9701}{mi}{65}%
}\hskip.1em\allowbreak{} 无重根，由第(1)问得\allowbreak{}\hskip.1em \RB{8.202}{8.294}{0.500}{%
\RG{0.000}{0.005}{10.9091}{mi}{65}%
}\hskip.1em\allowbreak{} 可对角化。

小结：无重根的极小多项式产生特征子空间的直和分解；共轭转置的交换关系排除了\allowbreak{}\hskip.1em \RB{59.399}{8.684}{3.236}{%
\RG{0.000}{0.002}{10.9091}{rm}{107}\RG{5.346}{0.002}{10.9091}{rm}{101}\RG{10.091}{0.002}{10.9091}{rm}{114}\RG{14.248}{0.002}{10.9091}{rm}{40}\RG{18.396}{0.002}{10.9091}{mi}{65}\RG{28.828}{0.002}{10.9091}{sy}{0}\RG{39.737}{0.002}{10.9091}{mi}{21}\RG{45.966}{0.002}{10.9091}{mi}{73}\RG{50.680}{-1.634}{7.9701}{mi}{110}\RG{55.823}{0.002}{10.9091}{rm}{41}%
}\hskip.1em\allowbreak{} 与\allowbreak{}\hskip.1em \RB{57.965}{8.684}{3.236}{%
\RG{0.000}{0.002}{10.9091}{rm}{73}\RG{3.851}{0.002}{10.9091}{rm}{109}\RG{12.797}{0.002}{10.9091}{rm}{40}\RG{16.962}{0.002}{10.9091}{mi}{65}\RG{27.394}{0.002}{10.9091}{sy}{0}\RG{38.303}{0.002}{10.9091}{mi}{21}\RG{44.532}{0.002}{10.9091}{mi}{73}\RG{49.246}{-1.634}{7.9701}{mi}{110}\RG{54.389}{0.002}{10.9091}{rm}{41}%
}\hskip.1em\allowbreak{} 的非零交向量。

% Original page 5.
\section*{2018S08}
设\allowbreak{}\hskip.1em \RB{8.715}{7.952}{0.728}{%
\RG{0.000}{0.001}{10.9091}{mi}{86}%
}\hskip.1em\allowbreak{} 是数域\allowbreak{}\hskip.1em \RB{10.035}{7.952}{0.500}{%
\RG{0.000}{0.001}{10.9091}{mi}{75}%
}\hskip.1em\allowbreak{} 上的\allowbreak{}\hskip.1em \RB{6.555}{5.312}{0.608}{%
\RG{0.000}{0.001}{10.9091}{mi}{110}%
}\hskip.1em\allowbreak{} 维线性空间，\allowbreak{}\hskip.1em \RB{7.122}{5.312}{2.855}{%
\RG{0.000}{0.001}{10.9091}{mi}{39}%
}\hskip.1em\allowbreak{} 是\allowbreak{}\hskip.1em \RB{8.715}{7.952}{0.728}{%
\RG{0.000}{0.001}{10.9091}{mi}{86}%
}\hskip.1em\allowbreak{} 上的线性变换，\allowbreak{}\hskip.1em \RB{49.060}{8.683}{3.237}{%
\RG{0.000}{0.001}{10.9091}{mi}{102}\RG{6.430}{0.001}{10.9091}{rm}{40}\RG{10.584}{0.001}{10.9091}{mi}{21}\RG{16.815}{0.001}{10.9091}{rm}{41}\RG{20.970}{0.001}{10.9091}{mi}{59}\RG{25.748}{0.001}{10.9091}{mi}{109}\RG{35.098}{0.001}{10.9091}{rm}{40}\RG{39.252}{0.001}{10.9091}{mi}{21}\RG{45.483}{0.001}{10.9091}{rm}{41}%
}\hskip.1em\allowbreak{} 分别是\allowbreak{}\hskip.1em \RB{7.122}{5.312}{2.855}{%
\RG{0.000}{0.001}{10.9091}{mi}{39}%
}\hskip.1em\allowbreak{} 的特征多项式和极小多项式。如果存在\allowbreak{}\hskip.1em \RB{8.715}{7.952}{0.728}{%
\RG{0.000}{0.001}{10.9091}{mi}{86}%
}\hskip.1em\allowbreak{} 的\allowbreak{}\hskip.1em \RB{7.122}{5.312}{2.855}{%
\RG{0.000}{0.001}{10.9091}{mi}{39}%
}\hskip.1em\allowbreak{} –不变子空间\allowbreak{}\hskip.1em \RB{25.527}{7.952}{2.605}{%
\RG{0.000}{0.001}{10.9091}{mi}{86}\RG{6.228}{-1.636}{7.9701}{rm}{49}\RG{10.546}{0.001}{10.9091}{mi}{59}\RG{15.324}{0.001}{10.9091}{mi}{86}\RG{21.558}{-1.636}{7.9701}{rm}{50}%
}\hskip.1em\allowbreak{} ，使得

\[
\RB{255.394}{8.074}{2.603}{%
\RG{0.000}{0.003}{10.9091}{mi}{86}\RG{11.632}{0.003}{10.9091}{rm}{61}\RG{22.971}{0.003}{10.9091}{mi}{86}\RG{35.940}{0.003}{10.9091}{sy}{8}\RG{46.849}{0.003}{10.9091}{mi}{86}\RG{57.395}{0.003}{10.9091}{mi}{59}\RG{83.999}{0.003}{10.9091}{rm}{100}\RG{89.933}{0.003}{10.9091}{rm}{105}\RG{92.901}{0.003}{10.9091}{rm}{109}\RG{103.673}{0.003}{10.9091}{mi}{86}\RG{117.249}{0.003}{10.9091}{mi}{60}\RG{128.588}{0.003}{10.9091}{rm}{100}\RG{134.522}{0.003}{10.9091}{rm}{105}\RG{137.490}{0.003}{10.9091}{rm}{109}\RG{148.261}{0.003}{10.9091}{mi}{86}\RG{155.080}{0.003}{10.9091}{mi}{59}\RG{181.686}{0.003}{10.9091}{rm}{100}\RG{187.621}{0.003}{10.9091}{rm}{105}\RG{190.588}{0.003}{10.9091}{rm}{109}\RG{201.361}{0.003}{10.9091}{mi}{86}\RG{214.936}{0.003}{10.9091}{mi}{60}\RG{226.274}{0.003}{10.9091}{rm}{100}\RG{232.209}{0.003}{10.9091}{rm}{105}\RG{235.176}{0.003}{10.9091}{rm}{109}\RG{245.949}{0.003}{10.9091}{mi}{86}\RG{252.767}{0.003}{10.9091}{mi}{59}\RG{29.198}{-1.634}{7.9701}{rm}{49}\RG{53.076}{-1.634}{7.9701}{rm}{50}\RG{109.900}{-1.634}{7.9701}{rm}{49}\RG{207.588}{-1.634}{7.9701}{rm}{50}%
}
\]
则称\allowbreak{}\hskip.1em \RB{8.715}{7.946}{0.734}{%
\RG{0.000}{-0.005}{10.9091}{mi}{86}%
}\hskip.1em\allowbreak{} 是\allowbreak{}\hskip.1em \RB{7.122}{5.306}{2.861}{%
\RG{0.000}{-0.005}{10.9091}{mi}{39}%
}\hskip.1em\allowbreak{} –可分解的，否则称\allowbreak{}\hskip.1em \RB{8.715}{7.946}{0.734}{%
\RG{0.000}{-0.005}{10.9091}{mi}{86}%
}\hskip.1em\allowbreak{} 是\allowbreak{}\hskip.1em \RB{7.122}{5.306}{2.861}{%
\RG{0.000}{-0.005}{10.9091}{mi}{39}%
}\hskip.1em\allowbreak{} –不可分解的。证明：\allowbreak{}\hskip.1em \RB{8.715}{7.946}{0.734}{%
\RG{0.000}{-0.005}{10.9091}{mi}{86}%
}\hskip.1em\allowbreak{} 是\allowbreak{}\hskip.1em \RB{7.122}{5.306}{2.861}{%
\RG{0.000}{-0.005}{10.9091}{mi}{39}%
}\hskip.1em\allowbreak{} –不可分解的充分必要条件是

\[
\RB{100.788}{10.056}{3.234}{%
\RG{93.489}{4.509}{7.9701}{mi}{107}\RG{0.000}{0.005}{10.9091}{mi}{102}\RG{6.430}{0.005}{10.9091}{rm}{40}\RG{10.583}{0.005}{10.9091}{mi}{21}\RG{16.815}{0.005}{10.9091}{rm}{41}\RG{23.993}{0.005}{10.9091}{rm}{61}\RG{35.338}{0.005}{10.9091}{mi}{109}\RG{44.681}{0.005}{10.9091}{rm}{40}\RG{48.835}{0.005}{10.9091}{mi}{21}\RG{55.066}{0.005}{10.9091}{rm}{41}\RG{62.245}{0.005}{10.9091}{rm}{61}\RG{73.589}{0.005}{10.9091}{mi}{112}\RG{78.950}{0.005}{10.9091}{rm}{40}\RG{83.104}{0.005}{10.9091}{mi}{21}\RG{89.335}{0.005}{10.9091}{rm}{41}\RG{98.161}{0.005}{10.9091}{mi}{59}%
}
\]
其中\allowbreak{}\hskip.1em \RB{19.323}{8.685}{3.235}{%
\RG{0.000}{0.003}{10.9091}{mi}{112}\RG{5.362}{0.003}{10.9091}{rm}{40}\RG{9.516}{0.003}{10.9091}{mi}{21}\RG{15.746}{0.003}{10.9091}{rm}{41}%
}\hskip.1em\allowbreak{} 是\allowbreak{}\hskip.1em \RB{10.035}{7.954}{0.500}{%
\RG{0.000}{0.003}{10.9091}{mi}{75}%
}\hskip.1em\allowbreak{} 上的首一不可约多项式，\allowbreak{}\hskip.1em \RB{25.440}{8.074}{1.992}{%
\RG{0.000}{0.003}{10.9091}{mi}{107}\RG{8.953}{0.003}{10.9091}{sy}{21}\RG{20.468}{0.003}{10.9091}{rm}{49}%
}\hskip.1em\allowbreak{} 。

证明：根据初等因子形式的循环分解定理，存在非零的\allowbreak{}\hskip.1em \RB{7.122}{5.312}{2.855}{%
\RG{0.000}{0.001}{10.9091}{mi}{39}%
}\hskip.1em\allowbreak{} –不变子空间\allowbreak{}\hskip.1em \RB{46.447}{7.952}{2.605}{%
\RG{0.000}{0.001}{10.9091}{mi}{85}\RG{7.254}{-1.635}{7.9701}{rm}{49}\RG{11.572}{0.001}{10.9091}{mi}{59}\RG{16.351}{0.001}{10.9091}{mi}{58}\RG{21.140}{0.001}{10.9091}{mi}{58}\RG{25.929}{0.001}{10.9091}{mi}{58}\RG{30.707}{0.001}{10.9091}{mi}{59}\RG{35.496}{0.001}{10.9091}{mi}{85}\RG{42.754}{-1.635}{7.9701}{mi}{115}%
}\hskip.1em\allowbreak{} ，使得

\[
\RB{87.847}{7.954}{2.603}{%
\RG{0.000}{0.003}{10.9091}{mi}{86}\RG{11.631}{0.003}{10.9091}{rm}{61}\RG{22.970}{0.003}{10.9091}{mi}{85}\RG{36.967}{0.003}{10.9091}{sy}{8}\RG{47.876}{0.003}{10.9091}{sy}{1}\RG{52.720}{0.003}{10.9091}{sy}{1}\RG{57.574}{0.003}{10.9091}{sy}{1}\RG{63.029}{0.003}{10.9091}{sy}{8}\RG{73.937}{0.003}{10.9091}{mi}{85}\RG{85.220}{0.003}{10.9091}{mi}{59}\RG{30.224}{-1.633}{7.9701}{rm}{49}\RG{81.192}{-1.633}{7.9701}{mi}{115}%
}
\]
并且每个限制变换\allowbreak{}\hskip.1em \RB{18.286}{8.677}{4.362}{%
\RG{0.000}{-0.005}{10.9091}{mi}{39}\RG{6.997}{-0.005}{10.9091}{sy}{106}\RG{10.028}{-1.641}{7.9701}{mi}{85}\RG{15.255}{-2.749}{5.9776}{mi}{106}%
}\hskip.1em\allowbreak{} 的特征多项式与极小多项式相同，均为\allowbreak{}\hskip.1em \RB{30.953}{9.501}{3.624}{%
\RG{0.000}{-0.005}{10.9091}{mi}{112}\RG{5.361}{-1.641}{7.9701}{mi}{106}\RG{9.360}{-0.005}{10.9091}{rm}{40}\RG{13.513}{-0.005}{10.9091}{mi}{21}\RG{19.745}{-0.005}{10.9091}{rm}{41}\RG{23.899}{3.954}{7.9701}{mi}{107}\RG{27.922}{2.845}{5.9776}{mi}{106}%
}\hskip.1em\allowbreak{} ，其中\allowbreak{}\hskip.1em \RB{8.865}{5.306}{3.624}{%
\RG{0.000}{-0.005}{10.9091}{mi}{112}\RG{5.361}{-1.641}{7.9701}{mi}{106}%
}\hskip.1em\allowbreak{} 首一不可约，\allowbreak{}\hskip.1em \RB{29.073}{8.075}{3.616}{%
\RG{0.000}{0.004}{10.9091}{mi}{107}\RG{5.555}{-1.632}{7.9701}{mi}{106}\RG{12.584}{0.004}{10.9091}{sy}{21}\RG{24.100}{0.004}{10.9091}{rm}{49}%
}\hskip.1em\allowbreak{} 。这一分解定理的出处见文末引用说明。

若\allowbreak{}\hskip.1em \RB{8.715}{7.952}{0.728}{%
\RG{0.000}{0.001}{10.9091}{mi}{86}%
}\hskip.1em\allowbreak{} 是\allowbreak{}\hskip.1em \RB{7.122}{5.312}{2.855}{%
\RG{0.000}{0.001}{10.9091}{mi}{39}%
}\hskip.1em\allowbreak{} –不可分解的，则\allowbreak{}\hskip.1em \RB{5.006}{5.312}{0.608}{%
\RG{0.000}{0.001}{10.9091}{mi}{115}%
}\hskip.1em\allowbreak{} 不能大于\allowbreak{}\hskip.1em \RB{4.973}{7.756}{0.500}{%
\RG{0.000}{0.001}{10.9091}{rm}{49}%
}\hskip.1em\allowbreak{} ；否则\allowbreak{}\hskip.1em \RB{11.006}{7.952}{2.135}{%
\RG{0.000}{0.001}{10.9091}{mi}{85}\RG{7.255}{-1.635}{7.9701}{rm}{49}%
}\hskip.1em\allowbreak{} 与\allowbreak{}\hskip.1em \RB{61.878}{7.952}{2.208}{%
\RG{0.000}{0.001}{10.9091}{mi}{85}\RG{7.254}{-1.635}{7.9701}{rm}{50}\RG{13.988}{0.001}{10.9091}{sy}{8}\RG{24.886}{0.001}{10.9091}{sy}{1}\RG{29.741}{0.001}{10.9091}{sy}{1}\RG{34.585}{0.001}{10.9091}{sy}{1}\RG{40.039}{0.001}{10.9091}{sy}{8}\RG{50.931}{0.001}{10.9091}{mi}{85}\RG{58.185}{-1.635}{7.9701}{mi}{115}%
}\hskip.1em\allowbreak{} 就给出两个非零不变子空间的直和分解：因此\allowbreak{}\hskip.1em \RB{24.370}{7.756}{0.608}{%
\RG{0.000}{0.001}{10.9091}{mi}{115}\RG{8.062}{0.001}{10.9091}{rm}{61}\RG{19.397}{0.001}{10.9091}{rm}{49}%
}\hskip.1em\allowbreak{} ，于是

\[
\RB{108.581}{10.054}{3.235}{%
\RG{97.808}{4.507}{7.9701}{mi}{107}\RG{101.831}{3.497}{5.9776}{rm}{49}\RG{0.000}{0.003}{10.9091}{mi}{102}\RG{6.430}{0.003}{10.9091}{rm}{40}\RG{10.584}{0.003}{10.9091}{mi}{21}\RG{16.815}{0.003}{10.9091}{rm}{41}\RG{23.993}{0.003}{10.9091}{rm}{61}\RG{35.338}{0.003}{10.9091}{mi}{109}\RG{44.681}{0.003}{10.9091}{rm}{40}\RG{48.835}{0.003}{10.9091}{mi}{21}\RG{55.066}{0.003}{10.9091}{rm}{41}\RG{62.245}{0.003}{10.9091}{rm}{61}\RG{73.590}{0.003}{10.9091}{mi}{112}\RG{83.269}{0.003}{10.9091}{rm}{40}\RG{87.422}{0.003}{10.9091}{mi}{21}\RG{93.654}{0.003}{10.9091}{rm}{41}\RG{106.074}{0.003}{10.9091}{mi}{58}\RG{78.951}{-1.633}{7.9701}{rm}{49}%
}
\]
反之，设\allowbreak{}\hskip.1em \RB{54.249}{9.509}{2.724}{%
\RG{0.000}{0.002}{10.9091}{mi}{102}\RG{9.461}{0.002}{10.9091}{rm}{61}\RG{20.799}{0.002}{10.9091}{mi}{109}\RG{33.173}{0.002}{10.9091}{rm}{61}\RG{44.511}{0.002}{10.9091}{mi}{112}\RG{49.872}{3.961}{7.9701}{mi}{107}%
}\hskip.1em\allowbreak{} 。若存在两个非零不变子空间，使\allowbreak{}\hskip.1em \RB{56.617}{7.953}{2.134}{%
\RG{0.000}{0.002}{10.9091}{mi}{86}\RG{11.632}{0.002}{10.9091}{rm}{61}\RG{22.970}{0.002}{10.9091}{mi}{86}\RG{29.198}{-1.634}{7.9701}{rm}{49}\RG{35.726}{0.002}{10.9091}{sy}{8}\RG{46.421}{0.002}{10.9091}{mi}{86}\RG{52.648}{-1.634}{7.9701}{rm}{50}%
}\hskip.1em\allowbreak{} ，记\allowbreak{}\hskip.1em \RB{17.122}{8.684}{3.298}{%
\RG{0.000}{0.002}{10.9091}{mi}{39}\RG{6.997}{0.002}{10.9091}{sy}{106}\RG{10.028}{-1.634}{7.9701}{mi}{86}\RG{14.527}{-2.743}{5.9776}{mi}{105}%
}\hskip.1em\allowbreak{} 的特征多项式为\allowbreak{}\hskip.1em \RB{8.071}{8.182}{2.724}{%
\RG{0.000}{0.002}{10.9091}{mi}{102}\RG{5.265}{-1.635}{7.9701}{mi}{105}%
}\hskip.1em\allowbreak{} 。选取与该直和分解相适应的基，其表示矩阵为分块对角阵，故

\[
\RB{66.875}{10.055}{2.722}{%
\RG{59.696}{4.508}{7.9701}{mi}{107}\RG{0.000}{0.004}{10.9091}{mi}{102}\RG{9.584}{0.004}{10.9091}{mi}{102}\RG{22.198}{0.004}{10.9091}{rm}{61}\RG{33.536}{0.004}{10.9091}{mi}{102}\RG{42.997}{0.004}{10.9091}{rm}{61}\RG{54.335}{0.004}{10.9091}{mi}{112}\RG{64.368}{0.004}{10.9091}{mi}{58}\RG{5.265}{-1.633}{7.9701}{rm}{49}\RG{14.849}{-1.633}{7.9701}{rm}{50}%
}
\]
因为\allowbreak{}\hskip.1em \RB{23.604}{8.176}{2.729}{%
\RG{0.000}{-0.004}{10.9091}{mi}{102}\RG{5.265}{-1.641}{7.9701}{rm}{49}\RG{9.584}{-0.004}{10.9091}{mi}{59}\RG{14.362}{-0.004}{10.9091}{mi}{102}\RG{19.635}{-1.641}{7.9701}{rm}{50}%
}\hskip.1em\allowbreak{} 均为正次数首一多项式，而\allowbreak{}\hskip.1em \RB{5.715}{5.307}{2.620}{%
\RG{0.000}{-0.004}{10.9091}{mi}{112}%
}\hskip.1em\allowbreak{} 不可约，所以存在正整数\allowbreak{}\hskip.1em \RB{24.287}{5.307}{2.609}{%
\RG{0.000}{-0.004}{10.9091}{mi}{97}\RG{5.607}{-1.641}{7.9701}{rm}{49}\RG{9.925}{-0.004}{10.9091}{mi}{59}\RG{14.704}{-0.004}{10.9091}{mi}{97}\RG{20.318}{-1.641}{7.9701}{rm}{50}%
}\hskip.1em\allowbreak{} ，满足

\[
\RB{184.325}{8.210}{2.728}{%
\RG{29.314}{4.502}{7.9701}{mi}{97}\RG{93.577}{4.502}{7.9701}{mi}{97}\RG{33.418}{3.492}{5.9776}{rm}{49}\RG{97.682}{3.492}{5.9776}{rm}{50}\RG{0.000}{-0.003}{10.9091}{mi}{102}\RG{12.614}{-0.003}{10.9091}{rm}{61}\RG{23.953}{-0.003}{10.9091}{mi}{112}\RG{37.660}{-0.003}{10.9091}{mi}{59}\RG{64.257}{-0.003}{10.9091}{mi}{102}\RG{76.878}{-0.003}{10.9091}{rm}{61}\RG{88.216}{-0.003}{10.9091}{mi}{112}\RG{101.925}{-0.003}{10.9091}{mi}{59}\RG{128.522}{-0.003}{10.9091}{mi}{97}\RG{140.878}{-0.003}{10.9091}{rm}{43}\RG{151.611}{-0.003}{10.9091}{mi}{97}\RG{164.567}{-0.003}{10.9091}{rm}{61}\RG{175.905}{-0.003}{10.9091}{mi}{107}\RG{181.818}{-0.003}{10.9091}{mi}{58}\RG{5.265}{-1.639}{7.9701}{rm}{49}\RG{69.529}{-1.639}{7.9701}{rm}{50}\RG{134.136}{-1.639}{7.9701}{rm}{49}\RG{157.218}{-1.639}{7.9701}{rm}{50}%
}
\]
对两个限制变换分别应用Cayley–Hamilton定理，有

\[
\RB{124.805}{11.602}{7.572}{%
\RG{28.111}{10.911}{9.0910}{ex}{12}\RG{28.111}{5.456}{9.0910}{ex}{12}\RG{20.667}{4.505}{7.9701}{mi}{97}\RG{24.771}{3.396}{5.9776}{mi}{105}\RG{0.000}{0.001}{10.9091}{mi}{112}\RG{5.362}{0.001}{10.9091}{rm}{40}\RG{9.516}{0.001}{10.9091}{mi}{39}\RG{16.513}{0.001}{10.9091}{rm}{41}\RG{28.111}{0.001}{9.0910}{ex}{12}\RG{42.012}{0.001}{10.9091}{rm}{61}\RG{53.347}{0.001}{10.9091}{rm}{48}\RG{80.511}{0.001}{10.9091}{rm}{40}\RG{84.664}{0.001}{10.9091}{mi}{105}\RG{91.338}{0.001}{10.9091}{rm}{61}\RG{102.673}{0.001}{10.9091}{rm}{49}\RG{108.018}{0.001}{10.9091}{mi}{59}\RG{112.803}{0.001}{10.9091}{rm}{50}\RG{118.149}{0.001}{10.9091}{rm}{41}\RG{122.298}{0.001}{10.9091}{mi}{58}\RG{31.141}{-5.909}{7.9701}{mi}{86}\RG{35.642}{-7.017}{5.9776}{mi}{105}%
}
\]
令\allowbreak{}\hskip.1em \RB{97.291}{8.683}{3.226}{%
\RG{0.000}{0.001}{10.9091}{mi}{97}\RG{8.978}{0.001}{10.9091}{rm}{61}\RG{20.662}{0.001}{10.9091}{rm}{109}\RG{29.564}{0.001}{10.9091}{rm}{97}\RG{34.910}{0.001}{10.9091}{rm}{120}\RG{40.725}{0.001}{10.9091}{sy}{102}\RG{46.179}{0.001}{10.9091}{mi}{97}\RG{51.786}{-1.636}{7.9701}{rm}{49}\RG{56.105}{0.001}{10.9091}{mi}{59}\RG{60.883}{0.001}{10.9091}{mi}{97}\RG{66.498}{-1.636}{7.9701}{rm}{50}\RG{70.816}{0.001}{10.9091}{sy}{103}\RG{79.641}{0.001}{10.9091}{mi}{60}\RG{91.325}{0.001}{10.9091}{mi}{107}%
}\hskip.1em\allowbreak{} ，则\allowbreak{}\hskip.1em \RB{25.043}{8.683}{3.237}{%
\RG{0.000}{0.001}{10.9091}{mi}{112}\RG{5.361}{0.001}{10.9091}{rm}{40}\RG{9.515}{0.001}{10.9091}{mi}{39}\RG{16.512}{0.001}{10.9091}{rm}{41}\RG{20.667}{3.959}{7.9701}{mi}{97}%
}\hskip.1em\allowbreak{} 在\allowbreak{}\hskip.1em \RB{25.527}{7.952}{2.605}{%
\RG{0.000}{0.001}{10.9091}{mi}{86}\RG{6.227}{-1.636}{7.9701}{rm}{49}\RG{10.545}{0.001}{10.9091}{mi}{59}\RG{15.323}{0.001}{10.9091}{mi}{86}\RG{21.558}{-1.636}{7.9701}{rm}{50}%
}\hskip.1em\allowbreak{} 上均为零，从而在\allowbreak{}\hskip.1em \RB{8.715}{7.952}{0.728}{%
\RG{0.000}{0.001}{10.9091}{mi}{86}%
}\hskip.1em\allowbreak{} 上为零。这与极小多项式为\allowbreak{}\hskip.1em \RB{9.738}{9.506}{2.616}{%
\RG{0.000}{-0.000}{10.9091}{mi}{112}\RG{5.362}{3.958}{7.9701}{mi}{107}%
}\hskip.1em\allowbreak{} 矛盾。因此\allowbreak{}\hskip.1em \RB{8.715}{7.951}{0.729}{%
\RG{0.000}{-0.000}{10.9091}{mi}{86}%
}\hskip.1em\allowbreak{} 是\allowbreak{}\hskip.1em \RB{7.122}{5.311}{2.856}{%
\RG{0.000}{-0.000}{10.9091}{mi}{39}%
}\hskip.1em\allowbreak{} –不可分解的。

小结：不可分解恰好对应一个不可约多项式幂的循环块；只有一个不可约因子并不足够，还必须使特征多项式与极小多项式相同。

% Original page 6.
\section*{2018S09}
设\allowbreak{}\hskip.1em \RB{8.715}{7.952}{0.728}{%
\RG{0.000}{0.001}{10.9091}{mi}{86}%
}\hskip.1em\allowbreak{} 是复数域上的\allowbreak{}\hskip.1em \RB{6.555}{5.312}{0.608}{%
\RG{0.000}{0.001}{10.9091}{mi}{110}%
}\hskip.1em\allowbreak{} 维线性空间，\allowbreak{}\hskip.1em \RB{7.122}{5.312}{2.855}{%
\RG{0.000}{0.001}{10.9091}{mi}{39}%
}\hskip.1em\allowbreak{} 是\allowbreak{}\hskip.1em \RB{8.715}{7.952}{0.728}{%
\RG{0.000}{0.001}{10.9091}{mi}{86}%
}\hskip.1em\allowbreak{} 上的线性变换，\allowbreak{}\hskip.1em \RB{8.562}{7.952}{0.728}{%
\RG{0.000}{0.001}{10.9091}{mi}{85}%
}\hskip.1em\allowbreak{} 是\allowbreak{}\hskip.1em \RB{8.715}{7.952}{0.728}{%
\RG{0.000}{0.001}{10.9091}{mi}{86}%
}\hskip.1em\allowbreak{} 的非零\allowbreak{}\hskip.1em \RB{7.122}{5.312}{2.855}{%
\RG{0.000}{0.001}{10.9091}{mi}{39}%
}\hskip.1em\allowbreak{} –不变子空间。设\allowbreak{}\hskip.1em \RB{10.288}{8.072}{2.289}{%
\RG{0.000}{0.001}{10.9091}{mi}{21}\RG{6.232}{-1.636}{7.9701}{rm}{48}%
}\hskip.1em\allowbreak{} 是限制变换\allowbreak{}\hskip.1em \RB{16.280}{8.683}{3.226}{%
\RG{0.000}{0.001}{10.9091}{mi}{39}\RG{6.997}{0.001}{10.9091}{sy}{106}\RG{10.028}{-1.636}{7.9701}{mi}{85}%
}\hskip.1em\allowbreak{} 的特征值，证明：\allowbreak{}\hskip.1em \RB{16.280}{8.683}{3.226}{%
\RG{0.000}{0.001}{10.9091}{mi}{39}\RG{6.998}{0.001}{10.9091}{sy}{106}\RG{10.028}{-1.636}{7.9701}{mi}{85}%
}\hskip.1em\allowbreak{} 的属于特征值\allowbreak{}\hskip.1em \RB{10.287}{8.072}{2.289}{%
\RG{0.000}{0.001}{10.9091}{mi}{21}\RG{6.231}{-1.636}{7.9701}{rm}{48}%
}\hskip.1em\allowbreak{} 的Jordan块的个数不超过\allowbreak{}\hskip.1em \RB{7.122}{5.311}{2.856}{%
\RG{0.000}{0.000}{10.9091}{mi}{39}%
}\hskip.1em\allowbreak{} 的属于特征值\allowbreak{}\hskip.1em \RB{10.287}{8.071}{2.289}{%
\RG{0.000}{0.000}{10.9091}{mi}{21}\RG{6.231}{-1.636}{7.9701}{rm}{48}%
}\hskip.1em\allowbreak{} 的Jordan块的个数。特别地，若\allowbreak{}\hskip.1em \RB{7.122}{5.311}{2.856}{%
\RG{0.000}{-0.000}{10.9091}{mi}{39}%
}\hskip.1em\allowbreak{} 的属于特征值\allowbreak{}\hskip.1em \RB{10.288}{8.071}{2.289}{%
\RG{0.000}{-0.000}{10.9091}{mi}{21}\RG{6.232}{-1.636}{7.9701}{rm}{48}%
}\hskip.1em\allowbreak{} 的Jordan块只有\allowbreak{}\hskip.1em \RB{4.973}{7.755}{0.500}{%
\RG{0.000}{0.000}{10.9091}{rm}{49}%
}\hskip.1em\allowbreak{} 个，那么\allowbreak{}\hskip.1em \RB{16.280}{8.682}{3.227}{%
\RG{0.000}{0.000}{10.9091}{mi}{39}\RG{6.998}{0.000}{10.9091}{sy}{106}\RG{10.028}{-1.637}{7.9701}{mi}{85}%
}\hskip.1em\allowbreak{} 的属于特征值\allowbreak{}\hskip.1em \RB{10.288}{8.071}{2.290}{%
\RG{0.000}{0.000}{10.9091}{mi}{21}\RG{6.232}{-1.637}{7.9701}{rm}{48}%
}\hskip.1em\allowbreak{} 的Jordan块也只有\allowbreak{}\hskip.1em \RB{4.973}{7.755}{0.500}{%
\RG{0.000}{0.000}{10.9091}{rm}{49}%
}\hskip.1em\allowbreak{} 个。

证明：对有限维复线性空间上的线性变换\allowbreak{}\hskip.1em \RB{8.049}{7.895}{0.502}{%
\RG{0.000}{-0.002}{10.9091}{mi}{84}%
}\hskip.1em\allowbreak{} ，记其属于\allowbreak{}\hskip.1em \RB{10.287}{8.069}{2.291}{%
\RG{0.000}{-0.002}{10.9091}{mi}{21}\RG{6.231}{-1.638}{7.9701}{rm}{48}%
}\hskip.1em\allowbreak{} 的Jordan块个数为\allowbreak{}\hskip.1em \RB{28.848}{8.680}{3.240}{%
\RG{0.000}{-0.002}{10.9091}{mi}{98}\RG{4.542}{-1.638}{7.9701}{mi}{84}\RG{10.569}{-0.002}{10.9091}{rm}{40}\RG{14.722}{-0.002}{10.9091}{mi}{21}\RG{20.954}{-1.638}{7.9701}{rm}{48}\RG{25.272}{-0.002}{10.9091}{rm}{41}%
}\hskip.1em\allowbreak{} 。在Jordan标准型中，\allowbreak{}\hskip.1em \RB{61.371}{8.679}{3.241}{%
\RG{0.000}{-0.003}{10.9091}{mi}{74}\RG{5.904}{-1.639}{7.9701}{mi}{114}\RG{10.060}{-0.003}{10.9091}{rm}{40}\RG{14.214}{-0.003}{10.9091}{mi}{21}\RG{20.445}{-1.639}{7.9701}{rm}{48}\RG{24.764}{-0.003}{10.9091}{rm}{41}\RG{31.342}{-0.003}{10.9091}{sy}{0}\RG{42.251}{-0.003}{10.9091}{mi}{21}\RG{48.482}{-1.639}{7.9701}{rm}{48}\RG{52.800}{-0.003}{10.9091}{mi}{73}\RG{57.511}{-1.639}{7.9701}{mi}{114}%
}\hskip.1em\allowbreak{} 的核为一维，而当\allowbreak{}\hskip.1em \RB{31.083}{8.308}{2.859}{%
\RG{0.000}{-0.003}{10.9091}{mi}{22}\RG{9.457}{-0.003}{10.9091}{sy}{54}\RG{9.457}{-0.003}{10.9091}{rm}{61}\RG{20.796}{-0.003}{10.9091}{mi}{21}\RG{27.027}{-1.639}{7.9701}{rm}{48}%
}\hskip.1em\allowbreak{} 时，\allowbreak{}\hskip.1em \RB{57.249}{8.679}{3.241}{%
\RG{0.000}{-0.003}{10.9091}{mi}{74}\RG{5.904}{-1.639}{7.9701}{mi}{114}\RG{10.060}{-0.003}{10.9091}{rm}{40}\RG{14.214}{-0.003}{10.9091}{mi}{22}\RG{20.641}{-0.003}{10.9091}{rm}{41}\RG{27.219}{-0.003}{10.9091}{sy}{0}\RG{38.128}{-0.003}{10.9091}{mi}{21}\RG{44.360}{-1.639}{7.9701}{rm}{48}\RG{48.678}{-0.003}{10.9091}{mi}{73}\RG{53.389}{-1.639}{7.9701}{mi}{114}%
}\hskip.1em\allowbreak{} 可逆，故

\[
\RB{125.722}{8.681}{3.239}{%
\RG{0.000}{-0.001}{10.9091}{mi}{98}\RG{10.568}{-0.001}{10.9091}{rm}{40}\RG{14.722}{-0.001}{10.9091}{mi}{21}\RG{25.271}{-0.001}{10.9091}{rm}{41}\RG{32.449}{-0.001}{10.9091}{rm}{61}\RG{43.795}{-0.001}{10.9091}{rm}{100}\RG{49.730}{-0.001}{10.9091}{rm}{105}\RG{52.697}{-0.001}{10.9091}{rm}{109}\RG{63.475}{-0.001}{10.9091}{rm}{107}\RG{68.821}{-0.001}{10.9091}{rm}{101}\RG{73.566}{-0.001}{10.9091}{rm}{114}\RG{77.723}{-0.001}{10.9091}{rm}{40}\RG{81.865}{-0.001}{10.9091}{mi}{84}\RG{92.033}{-0.001}{10.9091}{sy}{0}\RG{102.942}{-0.001}{10.9091}{mi}{21}\RG{113.492}{-0.001}{10.9091}{mi}{73}\RG{119.061}{-0.001}{10.9091}{rm}{41}\RG{123.215}{-0.001}{10.9091}{mi}{58}\RG{4.542}{-1.637}{7.9701}{mi}{84}\RG{20.953}{-1.637}{7.9701}{rm}{48}\RG{109.174}{-1.637}{7.9701}{rm}{48}%
}
\]
另一方面，限制变换的特征子空间满足

\[
\RB{255.651}{8.685}{3.235}{%
\RG{0.000}{0.003}{10.9091}{rm}{107}\RG{5.346}{0.003}{10.9091}{rm}{101}\RG{10.091}{0.003}{10.9091}{rm}{114}\RG{14.248}{0.003}{10.9091}{rm}{40}\RG{18.396}{0.003}{10.9091}{mi}{39}\RG{25.393}{0.003}{10.9091}{sy}{106}\RG{37.305}{0.003}{10.9091}{sy}{0}\RG{48.214}{0.003}{10.9091}{mi}{21}\RG{58.764}{0.003}{10.9091}{mi}{73}\RG{69.932}{0.003}{10.9091}{rm}{41}\RG{77.111}{0.003}{10.9091}{rm}{61}\RG{88.455}{0.003}{10.9091}{mi}{85}\RG{99.273}{0.003}{10.9091}{sy}{92}\RG{108.970}{0.003}{10.9091}{rm}{107}\RG{114.316}{0.003}{10.9091}{rm}{101}\RG{119.061}{0.003}{10.9091}{rm}{114}\RG{123.218}{0.003}{10.9091}{rm}{40}\RG{127.367}{0.003}{10.9091}{mi}{39}\RG{136.788}{0.003}{10.9091}{sy}{0}\RG{147.698}{0.003}{10.9091}{mi}{21}\RG{158.247}{0.003}{10.9091}{mi}{73}\RG{169.631}{0.003}{10.9091}{rm}{41}\RG{176.815}{0.003}{10.9091}{sy}{18}\RG{188.330}{0.003}{10.9091}{rm}{107}\RG{193.675}{0.003}{10.9091}{rm}{101}\RG{198.421}{0.003}{10.9091}{rm}{114}\RG{202.577}{0.003}{10.9091}{rm}{40}\RG{206.726}{0.003}{10.9091}{mi}{39}\RG{216.148}{0.003}{10.9091}{sy}{0}\RG{227.058}{0.003}{10.9091}{mi}{21}\RG{237.606}{0.003}{10.9091}{mi}{73}\RG{248.990}{0.003}{10.9091}{rm}{41}\RG{253.144}{0.003}{10.9091}{mi}{58}\RG{28.424}{-1.633}{7.9701}{mi}{85}\RG{54.446}{-1.633}{7.9701}{rm}{48}\RG{63.475}{-1.633}{7.9701}{mi}{85}\RG{153.929}{-1.633}{7.9701}{rm}{48}\RG{162.958}{-1.633}{7.9701}{mi}{86}\RG{233.288}{-1.633}{7.9701}{rm}{48}\RG{242.318}{-1.633}{7.9701}{mi}{86}%
}
\]
于是

\[
\RB{181.449}{25.907}{3.717}{%
\RG{0.000}{17.226}{10.9091}{mi}{98}\RG{17.793}{17.226}{10.9091}{rm}{40}\RG{21.947}{17.226}{10.9091}{mi}{21}\RG{32.496}{17.226}{10.9091}{rm}{41}\RG{39.675}{17.226}{10.9091}{rm}{61}\RG{51.020}{17.226}{10.9091}{rm}{100}\RG{56.955}{17.226}{10.9091}{rm}{105}\RG{59.922}{17.226}{10.9091}{rm}{109}\RG{70.690}{17.226}{10.9091}{rm}{107}\RG{76.046}{17.226}{10.9091}{rm}{101}\RG{80.792}{17.226}{10.9091}{rm}{114}\RG{84.948}{17.226}{10.9091}{rm}{40}\RG{89.090}{17.226}{10.9091}{mi}{39}\RG{96.087}{17.226}{10.9091}{sy}{106}\RG{107.999}{17.226}{10.9091}{sy}{0}\RG{118.908}{17.226}{10.9091}{mi}{21}\RG{129.457}{17.226}{10.9091}{mi}{73}\RG{140.626}{17.226}{10.9091}{rm}{41}\RG{28.177}{15.589}{7.9701}{rm}{48}\RG{99.117}{15.589}{7.9701}{mi}{85}\RG{125.139}{15.589}{7.9701}{rm}{48}\RG{134.169}{15.589}{7.9701}{mi}{85}\RG{4.542}{15.528}{7.9701}{mi}{39}\RG{9.578}{15.528}{7.9701}{sy}{106}\RG{11.724}{14.306}{5.9776}{mi}{85}\RG{39.681}{-0.002}{10.9091}{sy}{20}\RG{51.196}{-0.002}{10.9091}{rm}{100}\RG{57.130}{-0.002}{10.9091}{rm}{105}\RG{60.098}{-0.002}{10.9091}{rm}{109}\RG{70.865}{-0.002}{10.9091}{rm}{107}\RG{76.221}{-0.002}{10.9091}{rm}{101}\RG{80.967}{-0.002}{10.9091}{rm}{114}\RG{85.123}{-0.002}{10.9091}{rm}{40}\RG{89.267}{-0.002}{10.9091}{mi}{39}\RG{98.688}{-0.002}{10.9091}{sy}{0}\RG{109.597}{-0.002}{10.9091}{mi}{21}\RG{120.147}{-0.002}{10.9091}{mi}{73}\RG{131.530}{-0.002}{10.9091}{rm}{41}\RG{138.708}{-0.002}{10.9091}{rm}{61}\RG{150.053}{-0.002}{10.9091}{mi}{98}\RG{160.084}{-0.002}{10.9091}{rm}{40}\RG{164.239}{-0.002}{10.9091}{mi}{21}\RG{174.787}{-0.002}{10.9091}{rm}{41}\RG{178.942}{-0.002}{10.9091}{mi}{58}\RG{115.828}{-1.639}{7.9701}{rm}{48}\RG{124.858}{-1.639}{7.9701}{mi}{86}\RG{154.594}{-1.639}{7.9701}{mi}{39}\RG{170.469}{-1.639}{7.9701}{rm}{48}%
}
\]
若\allowbreak{}\hskip.1em \RB{48.232}{8.682}{3.715}{%
\RG{0.000}{-0.000}{10.9091}{mi}{98}\RG{4.542}{-1.637}{7.9701}{mi}{39}\RG{10.032}{-0.000}{10.9091}{rm}{40}\RG{14.186}{-0.000}{10.9091}{mi}{21}\RG{20.417}{-1.637}{7.9701}{rm}{48}\RG{24.735}{-0.000}{10.9091}{rm}{41}\RG{31.914}{-0.000}{10.9091}{rm}{61}\RG{43.259}{-0.000}{10.9091}{rm}{49}%
}\hskip.1em\allowbreak{} ，由\allowbreak{}\hskip.1em \RB{10.288}{8.071}{2.290}{%
\RG{0.000}{-0.000}{10.9091}{mi}{21}\RG{6.232}{-1.637}{7.9701}{rm}{48}%
}\hskip.1em\allowbreak{} 是\allowbreak{}\hskip.1em \RB{16.280}{8.682}{3.227}{%
\RG{0.000}{-0.000}{10.9091}{mi}{39}\RG{6.997}{-0.000}{10.9091}{sy}{106}\RG{10.028}{-1.637}{7.9701}{mi}{85}%
}\hskip.1em\allowbreak{} 的特征值，得\allowbreak{}\hskip.1em \RB{56.168}{8.682}{4.015}{%
\RG{0.000}{-0.000}{10.9091}{mi}{98}\RG{4.542}{-1.697}{7.9701}{mi}{39}\RG{9.578}{-1.697}{7.9701}{sy}{106}\RG{11.724}{-2.920}{5.9776}{mi}{85}\RG{17.793}{0.000}{10.9091}{rm}{40}\RG{21.947}{0.000}{10.9091}{mi}{21}\RG{28.178}{-1.637}{7.9701}{rm}{48}\RG{32.496}{0.000}{10.9091}{rm}{41}\RG{39.681}{0.000}{10.9091}{sy}{21}\RG{51.196}{0.000}{10.9091}{rm}{49}%
}\hskip.1em\allowbreak{} 。结合上式可知\allowbreak{}\hskip.1em \RB{55.992}{8.682}{4.015}{%
\RG{0.000}{0.000}{10.9091}{mi}{98}\RG{4.542}{-1.697}{7.9701}{mi}{39}\RG{9.577}{-1.697}{7.9701}{sy}{106}\RG{11.723}{-2.920}{5.9776}{mi}{85}\RG{17.792}{0.000}{10.9091}{rm}{40}\RG{21.947}{0.000}{10.9091}{mi}{21}\RG{28.177}{-1.637}{7.9701}{rm}{48}\RG{32.495}{0.000}{10.9091}{rm}{41}\RG{39.674}{0.000}{10.9091}{rm}{61}\RG{51.019}{0.000}{10.9091}{rm}{49}%
}\hskip.1em\allowbreak{} 。

小结：Jordan块的个数等于对应特征值的几何重数；限制到不变子空间后，特征子空间只能缩小。

引用说明

题面来源：谢启鸿《复旦大学高等代数习题集（2023版）》第1章§1.9，第60–61页。题中的高代教材[1]，按原书第206页参考文献，为谢启鸿、姚慕生、吴泉水《高等代数学（第四版）》，复旦大学出版社，2022。Cayley–Hamilton定理：可参见原书问题2016S06(4)（第25页），其伴随矩阵证明见第31–32页。初等因子形式的循环分解定理：其循环子空间分解表述见原书问题2020S09（第94页）；循环基与极小多项式的关系见问题2016S09(2)（第26页）。问题2018S08使用了这一标准结构定理。互素多项式的Bézout等式由欧几里得算法回代得到。无重根极小多项式判据及Jordan块个数与核维数的关系，已在本稿相应题目的证明中给出。


\end{document}
